% concatenated from FLArXive.tex
\documentclass[11pt]{amsart}
\usepackage{graphicx} % Required for inserting images
\usepackage[utf8]{inputenc}
\usepackage{amsmath}
\usepackage{xcolor}
\usepackage{amssymb}
\usepackage{amsmath}
\usepackage{amsfonts}
\usepackage{amstext}
\usepackage{amsthm}
\usepackage[english]{babel}
\usepackage[autostyle]{csquotes}
\usepackage[super]{nth}
\usepackage{hyperref}
\usepackage[margin=1.1in]{geometry}
\usepackage{txfonts}
%\usepackage[left]{showlabels}
\usepackage[normalem]{ulem}
\usepackage{parskip}

%\newcommand{\mycomment}[1]{}


\newtheorem{theorem}{Theorem}[section]
\newtheorem{proposition}[theorem]{Proposition}
\newtheorem{lemma}[theorem]{Lemma}
\newtheorem{corollary}[theorem]{Corollary}
\newtheorem{definition}[theorem]{Definition}
\newtheorem{remark}{Remark}
\newtheorem{claim}{Claim}
\newtheorem{example}{Example}[section]

\newenvironment{Assumptions}% Definition of assumptions
{%
\setcounter{enumi}{0}
\renewcommand{\theenumi}{(\textbf{A}.\arabic{enumi})}
\renewcommand{\labelenumi}{\theenumi}
\begin{enumerate}}%
{\end{enumerate} }

\makeatletter
\newcommand{\myitem}[1]{%
\medskip \item[#1]\protected@edef\@currentlabel{#1}%
} 

\newcommand{\rn}{\mathbb{R}^{n}}
\newcommand{\R}{\mathbb{R}}
\newcommand{\D}{\nabla}
\newcommand{\Fin}[1]{H^{p-1}(\nabla #1)\nabla H(\nabla #1)}
\newcommand{\fin}[1]{H^{p-1}( #1)\nabla H( #1)}
\newcommand{\wloc}[1]{W_{loc}^{1,p}(#1)}
\newcommand{\del}[1]{\delta(#1)}

 \numberwithin{equation}{section}

\providecommand*\lemmaautorefname{Lemma}
\providecommand*\propositionautorefname{Proposition}
\providecommand*\corollaryautorefname{Corollary}
\providecommand*\definitionautorefname{Definition}
\providecommand*\remarkautorefname{Remark}
\hypersetup{
    colorlinks=true,
    linkcolor= blue,
    citecolor=blue,
    pdfpagemode=FullScreen,
    }




\title[Boundary blowup solutions of Finsler p-Laplacian on bounded domains]{Boundary blowup solutions for the Finsler p-Laplacian: wellposedness and asymptotic behaviour}
%Large solutions for the Finsler p-Laplacian: existence, uniqueness, and asymptotic behavior near the boundary

%Large solutions of Finsler p-Laplacian on bounded domains

\author[N N Dattatreya]{N N Dattatreya}
\address{\parbox{.8\linewidth}
{{\textbf{N N Dattatreya}}\medskip \\
Indian Institute of Technology - Kanpur, India \medskip}}
\curraddr{}
\email{dattatreya21@iitk.ac.in}
\date{}


\begin{document}
\begin{abstract}
  We study the existence of \textit{large} (\textit{boundary blow-up}) solutions to semilinear equations involving the \textit{Finsler p-Laplacian} on bounded domains with sufficiently smooth boundaries. We establish a \textit{Keller–Osserman}-type condition that ensures the existence of such solutions, and show that this condition retains the same integrability as that of the p-Laplacian. We examine the influence of the anisotropic norm underlying the Finsler p-Laplacian on the boundary behaviour of the solution, then derive asymptotic estimates for large solutions near the boundary of the domain. Using these boundary asymptotics, we prove uniqueness results for power-type nonlinearities. 
\end{abstract}
\subjclass{35J \iffalse Elliptic equations and elliptic systems\fi, 35B44 \iffalse Blow-up in context of PDEs\fi, 35B40 \iffalse Asymptotic behavior of solutions to PDEs\fi,34A12 \iffalse Initial value problems, existence, uniqueness, continuous dependence and continuation of solutions to ordinary differential equations\fi, 35A02 \iffalse Uniqueness problems for PDEs: global uniqueness, local uniqueness, non-uniqueness\fi, 35A01 \iffalse Existence problems for PDEs: global existence, local existence, non-existence\fi. 
}
\keywords{Finsler p-Laplacian, Anisotropic operator, Existence, Nonexistence, Uniqueness, Boundary blow-up, Boundary asymptotic, Wulff shape, Anisotropic distance}
\maketitle
\section{Introduction}
We investigate positive large (boundary blow-up) solutions of semilinear equations driven by the Finsler p-Laplacian. Let $\Omega$ be a bounded domain in $\rn$. We consider positive solutions of the problem
\begin{equation}\label{finsler large eqn}
    \begin{cases}
        \text{div}\left(H^{p-1}(\nabla u)\nabla H(\nabla u)\right)=f(u)\quad &\text{in }\Omega\\
        u(x)\to\infty &\text{as dist}(x,\partial\Omega)\to 0.
    \end{cases}
\end{equation}
Where $2\leq p<\infty$ and $H\colon\R^n\rightarrow[0,+\infty)$ is a Minkowski norm (\cite{Shenbook2001}). i.e.,
\begin{itemize}
\item $H(x)>0$ for $x\ne 0$.
\item Positive homogeneity of degree 1: $H\left(tx\right)=|t|~H\left(x\right)$ for all $t\in \R$ and $x\in \R^{n}$.
    \item $H\in C^{2}\left(\R^n\setminus \{0\}\right)$.
     \item Strong convexity: The Hessian matrix $\nabla^{2}(H^{2}(x)/2)$ is positive definite for $x\in \rn\setminus\{0\}$.
    \end{itemize}
Further, the assumptions regarding nonlinearity are presented in subsection \ref{assumptions}.

The operator on the left-hand side of \eqref{finsler large eqn} is called the \textit{Finsler p-Laplacian}, denoted by $\Delta_{H}^{p}$.
 It is an anisotropic generalisation of the Laplace and the p-Laplace operators. In the case $p=2$, the operator $\Delta _{H}^{2}$ is called the Finsler Laplacian, which arise in Finsler geometry; we refer to \cite{BaoChernShen2000} for a comprehensive understanding of this geometry. Apart from geometric origin, anisotropic operators arise in various physical contexts, including minimal surface energy \cite{Taylor1978}, crystallography, crystal growth \cite{TaylorChanHandwerker1992}, image processing \cite{RudinOsherFatemi1992}, and the Willmore functional \cite{Ulrich2004}. The study of these anisotropic problems dates back to 1901, when Wulff \cite{Wulff1901} employed anisotropic geometry to investigate crystals. One of the basic objects of study in crystallography is \textit{Wulff shape}, which is the unit ball in the dual norm of $H$ and is a minimiser of an anisotropic functional defined on subsets of $\R^{n}$. For more in this direction, we refer to \cite{CozziFarinaValdinoci2014, Ulrich2004} and the references therein. The Wulff shape plays an important role in our analysis.

Large solutions to elliptic equations are studied extensively in various aspects, such as existence, uniqueness, and behaviour (asymptotics) of the solution and its gradient near the boundary of the domain. A necessary and sufficient condition for the existence of large solutions is the Keller-Osserman condition; for the Laplace equation, it reads as follows  
\begin{align*}
    \int^{\infty}_{r}\frac{ds}{\sqrt{F(s)}}<\infty,\quad \text{for all } r>0,
\end{align*}
where $F$ is the primitive of $f$ with $F(0)=0$.


The study of these problems traces back to L. Bieberbach’s 1916 work \cite{biberbach}, but the significant milestones were the independent contributions of Keller \cite{Keller1957} and Osserman \cite{Osserman1957}, both in 1957. The paper \cite{Keller1957} originated in electrodynamics \cite{osti_4157678} while \cite{Osserman1957} originated in geometry. A related development connecting large solutions to stochastic control can be found in \cite{lions}. For the existence, uniqueness, and asymptotics at the boundary of large solutions to semilinear equations, we refer to \cite{BandleMarcus1992, BandleMarcus1995, McKenna} and the references therein. For quasilinear equations involving the Laplace operator, we refer to \cite{BandleAntonioPorru1997, Lieberman}, and for the p-Laplacian, we refer to \cite{DattaIndro2026, diaz1993, materopaper} and the references therein.
While large solutions for the Laplacian and the p-Laplacian have been studied extensively, much less is known in the anisotropic framework associated with the Finsler operator; see \cite{DellaBlasio2017} for the Hamilton-Jacobi type equation involving the Finsler operator.

The present work extends classical results for semilinear equations to an anisotropic framework.

\subsection{Examples of Minkowski norms}
\begin{example}($2-$ norms)
    Consider
$H_{2}(x):=\left(\displaystyle\sum_{i=1}^{n}|x_{i}|^{2}\right)^{\frac{1}{2}}$.
\end{example}
\vspace{-4mm}
From this, one can obtain the Laplacian operator
       $ -div\left(\D u\right)=\Delta u=\Delta_{H_{2}}^{2}u$, and the $p$-Laplacian operator
    $-div\left(|\D u|^{p-2}\D u\right)=\Delta_{p}u=\Delta_{H_{2}}^{p}u$.
     \begin{example}
      Let $A$ be an $n\times n$ matrix. Take
      $H_{A}(x):=H_{2}(Ax)$.
  \end{example}
  From this, we can get a constant-coefficient elliptic operator of the form $- div\left(A\D u\right)=\Delta_{H_{A}}^{2}$.
  
   Before proceeding to the next example, notice that for $1\leq q<+\infty$, $q\neq 2$, the norms
    \begin{equation*}
H_{q}\left(x\right):=\left(\displaystyle\sum_{i=1}^{n}|x_{i}|^{q}\right)^{\frac{1}{q}}.
   \end{equation*} 
   fail to be Minkowski norms since $\nabla^{2}(H_{q}^{2}(x)/2)$ is only positive semi-definite and not positive definite, see \cite{BaoChernShen2000,Shenbook2001} for details. Consequently, such norms fall outside the framework considered in this article. Nevertheless, the admissible class of Minkowski norms $H$ satisfying the strong convexity condition remains quite rich.
  
   \begin{example}
       Let $\lambda\geq 0$ and $\mu>0$, consider 
       \begin{equation*}
           H_{\lambda,\mu}(x):=\sqrt{\lambda \sqrt{\sum_{i=1}^{n}|x_{i}|^{4}}+\mu\sum_{i=1}^{n}|x_{i}|^{2}}
       \end{equation*}
   \end{example}
Clearly, $H_{2}=H_{0,1}$. Moreover, for $\lambda_{1},\lambda_{2},\mu_{1},\mu_{2}>0$ such that $\lambda_{1}/\mu_{1}\neq \lambda_{2}/\mu_{2}$, the norms $H_{\lambda_{1},\mu_{1}}$ and $H_{\lambda_{2},\mu_{2}}$ are non-isomorphic, see \cite{MezeiVas2019}.
\begin{example}
    Let $n=m_{1}+\cdots +m_{r}$ and $1\leq q,p_{i},\lambda_{i}<+\infty$ for $1\leq i\leq r$, denote $P=\left(p_{1},\cdots,p_{r}\right)$ and $\Lambda=\left(\lambda_{1},\cdots,\lambda_{r}\right)$. Consider 
    \begin{equation*}
        H=H_{q,P,\Lambda}:=\left(\displaystyle \sum_{i=1}^{r}\lambda_{i}\left(\sum_{j=s_{i-1}}^{m_{i}}|x_{j}|^{p_{i}}\right)^{\frac{q}{p_{i}}}\right)^{\frac{1}{q}},
    \end{equation*}
   where $s_{0}=0$ and $s_{k}=\displaystyle\sum_{i=1}^{k}m_{i}$.
\end{example}
\begin{example}(Randers norm)\cite[Example 1.2.2]{Shenbook2001}
    Let $T$ be a linear functional on $\rn$. Define
    \begin{equation*}
        H(x)=H_{2}(x)+Tx.
    \end{equation*}
    Then $H$ is a Minkowski norm if and only if $\|T\|\leq 1$.
\end{example}


\subsection{Assumptions}\label{assumptions}
Throughout the paper, we assume that $p\geq 2$ and that $\Omega$ is a bounded domain in $\R^{n}$ with $C^{2}$ boundary.

We assume the following on the nonlinearity:
\begin{enumerate}
    \myitem {$\textbf{(F)}$}\label{a1}
    \textit{$f:[0,\infty)\to [0,\infty)$ is continuous and strictly increasing, with $f(0)=0$ and $f(r)\to \infty$ as $r\to\infty$.}
\end{enumerate}
In the classical spirit, we assume the Keller-Osserman condition associated with the Finsler p-Laplace equation: Define the primitive of $f$ by
\begin{equation}\label{eqn primitive of f}
    F(x)=\int_{0}^{x}f(s)\ ds.
\end{equation}
Then we assume
%\vspace{-0.2cm}
\begin{enumerate}
    \myitem{$\textbf{(KO)}$} \label{KO}$ \Psi_{H,p}(r)=\left(\frac{p-1}{p}\right)^{1/p}\int_{r}^{\infty}\frac{ds}{\left\{F(s)\right\}^{1/p}}<\infty, \quad \text{for all }~ r>0$.
\end{enumerate}
\begin{remark}
    Remarkably, despite the anisotropic nature of the operator, the Keller–Osserman condition retains the same integrability structure as in the p-Laplacian case. This is because all anisotropic norms are equivalent to the Euclidean norm. The above Keller–Osserman condition coincides with the Keller–Osserman condition for the p-Laplacian; see \cite{DattaIndro2026}. See also Remark \ref{remark 1}, and proof of Theorem \ref{existence on annulus} or Proposition \ref{prop annuls asym} in Section \ref{asym}. 
    \end{remark}
    \begin{example}\label{example f=tq}
         If $f(t)=t^{q}$, then \ref{KO} implies $q>p-1$. And, 
    \begin{equation*}
        \Psi_{H,p}(t)= \left(\frac{(p-1)(q+1)}{p}\right)^{1/p}\frac{p}{q+1-p}\frac{1}{t^{\frac{q+1-p}{p}}}.
    \end{equation*}
    \end{example}
We further consider the two possibilities about the integrability of $\Psi_{H,p}$ near $0$:
\begin{enumerate}
        \myitem{$\textbf{(A1)}$}\label{i1} $\int_{0^{+}}\frac{ds}{\{F(s)\}^{1/p}}=\infty$ ($Osgood~ Condition$).
    \myitem{$\textbf{(A2)}$}\label{i2} $\int_{0^{+}}\frac{ds}{\{F(s)\}^{1/p}}<\infty$.
\end{enumerate}

\subsection{Main results}
Consider the following notion of solution.
\begin{definition}\label{dfn solution}
    A function $u\in W_{loc}^{1,p}(\Omega)$ is said to solve 
    \begin{equation}\label{finsler eqn}
         \text{div}\left(H^{p-1}(\nabla u)\nabla H(\nabla u)\right)=f(u)\quad \text{in }\Omega,
    \end{equation}
    if for every pre-compact set $\Omega'\subset\subset \Omega$ 
    \begin{equation*}
        -\int_{\Omega'}\Fin{u}\cdot \D \phi\ dx=\int_{\Omega'}f(u)\phi\ dx \quad \text{for all }\phi\in W_{0}^{1,p}(\Omega').
    \end{equation*}
    provided $f(u)\in L_{loc}^{p'}(\Omega)$.
\end{definition}
\begin{remark}
    For a fixed $u\in W_{loc}^{1,p}(\Omega)$, the mapping $v\mapsto \int_{\Omega}\Fin{u}\cdot \nabla v $ for any $v\in W^{1,p'}(\Omega)$ is bounded and linear   due to Lemma \ref{lemma_properties of H} below.
\end{remark}

Comparison principle plays an important role in proving both existence and asymptotic results. We prove a comparison principle 

\begin{theorem}\label{comparison principle}
    Let $\Omega$ be a bounded domain in $\rn$ with $C^{2}$ boundary, $p\geq 2$, and $u,v\in W_{loc}^{1,p}(\Omega)\cap C(\Omega)$ such that 
    \begin{equation}\label{ineq for comp princ}
       - \text{div}\left(\Fin{u}\right)+f(u)\leq -\text{div}\left(\Fin{v}\right)+f(v) \quad\text{weakly}.
    \end{equation}
  Further, suppose that $u,v$ satisfy 
    \begin{equation}\label{boundary condition for comp princ}
        \limsup_{x\to z}\frac{u(x)}{v(x)}\leq 1,\quad \text{uniformly for all }z\in \partial\Omega.
    \end{equation}
    Then $u\leq v$ in $\Omega$.
    \end{theorem}
 We refer to Definition \ref{dfn solution} for the weak notion. Further, we prove the following existence result
  \begin{theorem}\label{main theorem existence}
     Let $\Omega$ be a bounded domain in $\rn$ with $C^{2}$ boundary. Assume \ref{a1}
 and \ref{KO} then \eqref{finsler large eqn} admits a solution in $W^{1,p}_{loc}(\Omega)$.
 \end{theorem}
 Next, we prove the boundary asymptotic for any large solutions. Before stating the result, the dual norm $H_{0}$ of $H$ is 
\begin{equation}\label{dual dfn}
    H_{0}\left(\xi\right):=\displaystyle\sup_{x\neq 0}\frac{\left<\xi, x\right>}{H\left(x\right)} \quad \text{for all } \xi\in \R^{n}.
\end{equation}
 \begin{theorem}\label{dual asym}
   Let $u$ be a solution of \eqref{finsler large eqn} on a bounded domain $\Omega$ with $C^{2}$ boundary. Let $\delta_{H_{0}}(x):=\min\{H_{0}(x-z)~|~z\in \partial\Omega\}$ %and consider the normalized $\Psi_{H,p}$,
   %\begin{equation}\label{eqn tilda psi}
    %    \Psi_{H,p}(r)=\left(\frac{p-1}{p}\right)^{1/p}\int_{r}^{\infty}\frac{ds}{\left\{F(s)\right\}^{1/p}}=\frac{1}{\gamma} \Psi_{H,p}(r).
  %  \end{equation}
    Then
       \begin{equation*}
        \lim_{\delta_{H_{0}}(x)\to 0}\frac{\Psi_{H,p}(u(x))}{\delta_{H_{0}}(x)}= 1.
    \end{equation*}   
 \end{theorem}
 Finally, defining $\Phi:=\Psi^{-1}$, we prove the following boundary blow-up result.
  \begin{corollary}\label{thrm asym for f=tq}
        Let $\Omega$ be a domain with $C^{2}$ boundary, and let $u$ be a solution of \eqref{finsler large eqn} with $f(t)=t^{q}$, $q>p-1$. Then 
        \begin{equation}\label{asym f=tq}
            \lim_{\delta_{H_{0}}(x)\to 0}\frac{u(x)}{\Phi(\delta_{H_{0}}(x))}=1.
        \end{equation}
    \end{corollary}
    As a consequence, we obtain the following uniqueness result.
 \begin{theorem}\label{thrm uniqueness}
     Let $\Omega$ be a bounded domain with $C^{2}$ boundary, then the equation \eqref{finsler large eqn} with $f(t)=t^{q}$, $q>p-1$ admits exactly one solution.
 \end{theorem}

\subsection{Outline}
 The existence of solutions to the finite Dirichlet problem, obtained via standard minimisation techniques, together with a brief discussion of regularity, is presented in the Appendix, since they are not the main focus of the paper.

The proof of the existence theorem uses the Browder–Minty-type argument. The proof of the comparison principle follows a standard argument; however, for completeness, we provide a full proof, since, to the best of our knowledge, this version of the comparison principle is not available in the literature.
The proof of the local bound in Section \ref{higher dim} follows a similar approach to \cite{DattaIndro2026}, but our proof is a somewhat simplified with a slightly different choice of super solution. Moreover, since the symmetry of the solution is with respect to the underlying anisotropic norm, we obtain a local bound using Wulff balls. The proof of asymptotic behaviour relies heavily on the distance defined using the Minkowski norm, the Wulff shape, and the geometry of the domain in Minkowski space. Precisely, we need to compare the solution with the distance from the boundary with respect to $H_{0}$ norm. While doing so, one needs to use uniform ball conditions associated with the boundary of a $C^{2}$ domain in the Minkowski space; see the first two paragraphs in Section \ref{asym}.

The focus of this work is on extending the theory of large solutions to anisotropic operators of Finsler p-Laplacian type. While the Keller–Osserman condition remains unchanged, the analysis requires handling anisotropic geometry induced by the norm $H$ and its dual $H_{0}$. For example, the solution exhibits radial symmetry on a Wulff ball associated with $H_{0}$. Consequently, the barrier construction is on a similar $H_{0}$ Wulff ball of radius $r$. Similarly, the boundary asymptotics are reformulated in terms of the distance from the boundary induced by $H_{0}$. This is where the geometry of the Wulff shape, and the distance from a closed set in Minkowski space (\cite{CrastaMalusa2007}) play an important role.

   In Section \ref{prilms}, we present preliminary results concerning the Minkowski norm $H$ and its dual $H_{0}$, and establish the strict monotonicity of the Finsler p-Laplace operator. In Section \ref{1d}, we derive condition \ref{KO} in the one-dimensional setting and prove Theorem \ref{main theorem existence} for $n=1$; see also \cite{DattaIndro2026} for a similar construction for a second-order quasilinear equation. The proof of this existence theorem in one dimension follows the same approach as in \cite{DattaIndro2026} except some technical differences. In the same section, we prove Theorem \ref{dual asym} for $n=1$.
Subsequently, in Section \ref{higher dim}, we establish the comparison principle (see Theorem \ref{comparison principle}), derive estimates for the solution and its gradient, and prove the existence result stated in Theorem \ref{main theorem existence}. In Section \ref{asym}, we give the proof of Theorem \ref{dual asym} for $n\geq2$ and discuss how the anisotropy underlying the Finsler operator plays a vital role in the theory. In Section \ref{uniqueness}, we use the asymptotic results to prove Corollary  \ref{thrm asym for f=tq} and finally, the uniqueness result, Theorem \ref{thrm uniqueness}.



\section{Preliminaries}\label{prilms}
This section is devoted to the discussion related to the properties of $H$ and $H_{0}$ and the monotonicity of the Finsler p-Laplace operator. First, we fix some notation.


 \textbf{Notation:}
Let $\Omega$ be a bounded domain
\begin{itemize}
\item We denote $H_{0}$ to be the dual norm of $H$ as defined in \eqref{dual dfn}.
    \item The distance from the boundary of $\Omega$ with respect to the Euclidean norm is $\delta(x):=$dist$(x,\partial\Omega)$. And with respect to the dual norm $H_{0}$ is $\delta_{H_{0}}(x):=\inf\{H_{0}(x-z)~|~z\in \partial\Omega\}$.
    \item We denote $\theta_{1}=\min_{|x|=1}H(x)$ and $\theta_{2}=\max_{|x|=1}H(x)$.
    \item The ball with respect to $H_{0}$ is $\mathcal{W}_{r}(x_{0}):=\{x\in \R^{n}~|~ H_{0}(x-x_{0})<r\}$ for any $x_{0}\in \R^{n}$. In the literature $\mathcal{W}_{1}$ is known as \textit{Wulff Shape}. 
    \item For any $\delta>0$, we denote $\Omega_{\delta}:=\{x\in \Omega~|~\delta_{H_{0}}(x)<\delta\}$.
    \item By $A\subset\subset B$, we mean that $A$ is a pre-compact subset of $B$.
\end{itemize}

The following properties hold for $H$ and $H_{0}$.

 \begin{lemma}\label{lemma_properties of H}
    Let $H$ be a Minkowski norm. Then
    \begin{enumerate}
        \item For $\theta_{1}$ and $\theta_{2}$ as in the Notation, we have $\theta_{1}|z|\leq H\left(z\right)\leq \theta_{2}|z|$ and $0<\theta_{1}\leq \theta_{2}$.
        \item $H\left(x+y\right)\leq H\left(x\right)+H\left(y\right)$.  
\item There exists a constant $c>0$ such that $|\D H|\leq c$.
\item For any vector $x$, $\left<\D H(x), x\right>=H\left(x\right)$.
\end{enumerate}
\end{lemma}
\begin{proof}
    From the homogeneity,  $H\left(z\right)=|z|H\left(\frac{z}{|z|}\right)$, the first one follows from the definitions of $\theta_{1}$ and $\theta_{2}$. Strict inequality $0<\theta_{1}$ follows from the continuity, strict positivity, and homogeneity of $H$. In particular, when $\theta_{1}=\theta_{2}$, $H=\theta_{1}H_{2}$.
The second one holds since strong convexity implies convexity.
    Further, we have by $(2)$ and the homogeneity 
    \begin{equation*}
    \frac{H\left(x+he_{i}\right)-H\left(x\right)}{h}\leq H\left(e_{i}\right). 
    \end{equation*}
    Therefore, the third one holds thus with $c=\left(\displaystyle\sum_{i=1}^{n}H^{2}(e_{i})\right)^{\frac{1}{2}}$.
     Property $(4)$ can be obtained by differentiating $H\left(tx\right)=t H\left(x\right), t>0$, with respect to $t$ and substituting $t=1$.
\end{proof}

\begin{lemma}\label{H0} $H$ and $H_{0}$ interact in the following way
    \begin{enumerate}
     \item $H$-H\"older inequality: $\left<\xi, x\right>\leq H_{0}\left(\xi\right)H\left(x\right)$.
     \item $H_{0}$ is a Minkowski norm.
        \item Dual of $H_{0}$ is $H$.
        \item $H_{0}\left(\D H(x)\right)=1$ and $H\left(\D H_{0}(x)\right)=1$ for any $x$.
        \item $H_{0}(x)\nabla H(\nabla H_{0}(x))=x$ for any $x$.
    \end{enumerate}
\end{lemma}
\begin{proof}
    If $x=0$, then $H$-H\"older holds trivially. Otherwise, it holds by 
    \begin{equation*}
        \left<\xi, x\right>=\frac{\left<\xi , x\right>H\left(x\right)}{H\left(x\right)},
    \end{equation*}
    and the definition of $H_{0}$.
    For the proof of the remaining properties, we refer to \cite[Proposition 1.3]{DissertationXia}, \cite[Section 2.1]{BePo1996}, and the references therein.
\end{proof}

Strict monotonicity is one of the key aspects in the proof of the existence of a solution to \eqref{finsler large eqn}. We dedicate the rest of the results to proving that the operator we have considered is strictly monotonic.
\begin{lemma}\label{euler inequality}
    Let $x,y\in \rn\setminus\{0\}$. Then we have 
    \begin{equation}\label{convixity and tangent plane}
        H(x)\geq \left<\D H(y), x\right>
    \end{equation}
\end{lemma}
\begin{proof}
     Strong convexity of $H$ implies convexity. Thus, for any $t\in (0,1)$,
    \begin{equation*}
        H(x)-H(y)\geq \frac{H(y+t(x-y))-H(y)}{t}.
    \end{equation*}
       As directional derivatives exist, by taking $t\to 0$, we obtain  
     \begin{equation}\label{eqn hyperplane}
         H(x)-H(y)\geq \left<\D H(y), x-y\right>.
     \end{equation}
     Using the last implication of Lemma \ref{lemma_properties of H} on $H(y)$, we obtain  \eqref{convixity and tangent plane}.
\end{proof}
\begin{lemma}\label{lemma monotone operator}
    For any $x,y \in \rn$
    \begin{equation*}
        \left< \fin{x}-\fin{y}, x-y\right>\geq 0.
    \end{equation*}
\end{lemma}
\begin{proof}
    If $y=0$, then by Lemma \ref{lemma_properties of H} the inequality holds in view of the positivity of $H$. The case $x=0$ is handled similarly. Therefore, we may assume that $x\neq 0$ and $y\neq 0$. By $(4)$ in Lemma \ref{lemma_properties of H} and Lemma \ref{euler inequality} 
\begin{align*}
    &\left<\fin{x}-\fin{y}, x-y\right>\\
    &= H^{p}(x)-\left<\fin{x},  y\right>
    +H^{p}(y)-\left<\fin{y}, x\right>\\
    &\geq H^{p}(x)-H^{p-1}(x) H(y)
    +H^{p}(y)-H^{p-1}(y) H(x)\\
    &\geq \left(H(x)-H(y)\right)\left(H^{p-1}(x)-H^{p-1}(y)\right)\geq 0. \qedhere
\end{align*}
 \end{proof}
 Notice that Lemma \ref{lemma monotone operator} holds if $H$ is only assumed to be convex. Strict convexity implies strict monotonicity, as implied in the next three results. We first consider an improvement of Lemma \ref{euler inequality}, which can be found in \cite[Proposition 4.6]{Ohta2009}.
\begin{lemma}
    There exists $C\geq 1$ such that for all $x,y\in\rn $ 
    \begin{equation*}
        \frac{H^{2}(x)+H^{2}(y)}{2}\geq H^{2}\left(\frac{x+y}{2}\right)+\frac{1}{4C^{2}}H^{2}(x-y).
    \end{equation*}
\end{lemma}
The following two lemmas are from \cite{BalGarainMukarjee2021}.
\begin{lemma}
    Let $p\geq 2$, there exists a constant $c>0$ such that for all $x,y\in\rn$
    \begin{equation}\label{eqn improved hyperplane}
        H^{p}(x)\geq H^{p}(y)+\left<\nabla H^{p}(y)\cdot x-y\right>+cH^{p}(x-y).
    \end{equation}
\end{lemma}
\begin{proof}
    Since $p\geq 2$, we use the inequality $a^{p}+b^{p}\leq (a^{2}+b^{2})^{p/2}$ for $a,b\geq 0$ , homogeneity of $H$, previous lemma and $(a+b)^{p}\leq 2^{p-1}(a^{p}+b^{p})$ to obtain
    \begin{equation*}
    \begin{split}
       H^{p}\left(\frac{x+y}{2}\right)+\frac{1}{C^{p}}H^{p}\left(\frac{x-y}{2}\right)&\leq \left\{H^{2}\left(\frac{x+y}{2}\right)+\frac{1}{C^{2}}H^{2}\left(\frac{x-y}{2}\right)\right\}^{\frac{p}{2}} \\
       &= \left\{H^{2}\left(\frac{x+y}{2}\right)+\frac{1}{4C^{2}}H^{2}\left(x-y\right)\right\}^{\frac{p}{2}}\\
       &\leq \frac{1}{2^{\frac{p}{2}}}\left\{H^{2}(x)+H^{2}(y)\right\}^{\frac{p}{2}}\\
       &\leq \frac{H^{p}(x)+H^{p}(y)}{2}.
    \end{split}
    \end{equation*}
    Since $H^{p}$ convex, same logic as in Lemma \ref{euler inequality} and replacing $x$ by $(x+y)/2$ in \eqref{eqn hyperplane} with $H^{p}$ instead of $H$, and using it in the above inequality we obtain 
    \begin{equation*}
        H^{p}(y)+\frac{1}{2}\left<\nabla H^{p}(y), x-y\right>+\frac{1}{C^{p}}H^{p}\left(\frac{x-y}{2}\right)\leq \frac{H^{p}(x)+H^{p}(y)}{2}.
    \end{equation*}
    The lemma follows from the homogeneity, with $c=1/2C^{2}$.
\end{proof}
\begin{corollary}\label{strict convexity of H}
    Let $p\geq 2$. There exists a positive constant $c=c(\theta_{1},\theta_{2},p)$ such that for any $x,y\in \rn$
    \begin{equation*}
          \left<\fin{x}-\fin{y}, x-y\right>\geq c H^{p}(x-y).
    \end{equation*}
    \end{corollary}
\begin{proof}
Interchanging $x$ and $y$ in \eqref{eqn improved hyperplane} and adding the resulting inequality to \eqref{eqn improved hyperplane}, the corollary follows.
\end{proof}
\section{Analysis in one dimension}\label{1d}
In this section, we restrict ourselves to one dimension and show that the Keller–Osserman condition is necessary and sufficient for the existence of large solutions on an interval, consequently proving Theorem \ref{main theorem existence} for $n=1$. Further, we explore the behaviour of a solution near the endpoints of the interval and subsequently prove Theorem \ref{dual asym} for $n=1$.

\subsection{Existence:}
Let $a,b\in \R$. We consider the problem 
\begin{equation}\label{eqn in one dim without boundary}
    \left(H^{p-1}(u')H'(u')\right)'=f(u)\quad \text{on } (a,b).
\end{equation}
In dimension one, any Minkowski norm, being positively 1-homogeneous, must be of the form $H(t)=\gamma |t|$ for some $\gamma>0$. The above equation reduces to  
\begin{equation}\label{eqn reduced in one dim}
    \gamma^{p}\left(|u'|^{p-2}u'\right)'=f(u)\quad \text{on }(a,b).
\end{equation}
If $u$ is a classical solution,
since $f\geq 0$ by \ref{a1}, the equation implies $(p-1)\gamma^{p}|u'|^{p-2}u''\geq 0$. That is,  $u$ must be convex. Therefore, we look for a convex solution.

    With this at hand, we prove Theorem \ref{main theorem existence} for $n=1$.
    \begin{theorem}\label{thrm existence in one dim with osgood}
        Assume \ref{KO}, \ref{a1}, and \ref{i1}. Then  
        \begin{equation}\label{eqn in one dim}
        \begin{cases}
      \left(H^{p-1}(v')H'(v')\right)'=f(v)\quad \text{on } (a,b). \\
      v(x)\to \infty \quad \text{as } x\to a \text{ or }b
        \end{cases}
        \end{equation}
        admits a solution.
    \end{theorem}
    \begin{proof}Let $c_{m}$ be the point at which $u$ attains its minimum. Denote $v_{0}=u(c_{m})$. We prove the existence in three steps; in the first step, we associate a function $v_{0} \mapsto \ell(v_{0})$ which determines an interval $\left(-\ell(v_{0}),\ell(v_{0})\right)$ where unique solution exists. In the next step we show that the map $\ell:(0,\infty)\to (0,\infty)$ is a bijection, which implies that given any $\ell>0$, we can construct the solution on $(-\ell,\ell)$. Finally, in the last step we use translations to get unique solution on any interval $(a,b)$.

    \textit{Step 1:} We consider the second-order initial value problem
\begin{equation}\label{Ode1}
 \begin{cases}
     \gamma^{p}\left(|v'|^{p-2}v'\right)'=f(v)\quad \text{for all } x\geq c_{m}\\
     v'(c_{m})=0 \quad \text{and} \quad v(c_{m})=v_{0}.
 \end{cases}
 \end{equation}
 Multiplying the above equation by $v'$, integrating from $c_{m}$ to $x$, and using the change of variables $s=v(t)$, we obtain 
 \begin{equation*}
     (p-1)\gamma^{p}\int_{0}^{v'(x)}t^{p-1}\ dt=F\left(v(t)\right)-F(v_{0})
 \end{equation*}
 where $F$ is defined in \eqref{eqn primitive of f}.
Rewriting the above equation
 \begin{equation}\label{eqn quotient}
     \left(\frac{p-1}{p}\right)^{1/p}\frac{\gamma v'(x)}{\left\{F\left(v(x)\right)-F(v_{0})\right\}^{1/p}}=1.
 \end{equation}
 Integrating from $c_{m}$ to $x$ and using a change of variables,
 \begin{equation*}
     \left(\frac{p-1}{p}\right)^{1/p}\gamma \int_{v_{0}}^{v(x)}\frac{ds}{\left\{F(s)-F(v_{0})\right\}^{1/p}}=x-c_{m}.
 \end{equation*}
 \begin{remark}
     If $v(x)\to +\infty$ as $x\to b$, then by the right hand side \ref{KO} holds. On the other hand if \ref{KO} holds, then $v$ has to blow up at a point in $\R$ since otherwise right hand side can tend to $+\infty$.
 \end{remark}
  \begin{remark}\label{remark 1}
     The appearance of anisotropy in the Keller-Osserman condition occurs only through the multiplicative constant $\gamma$; consequently, the integrability requirement is invariant under the choice of anisotropy $H$.
 \end{remark}
Uniqueness of $v$ follows by the positivity of the integrand. %This implies that $v(x)=u(x)$ for $x\geq c_{m}$. 

Let $(c_{m},\ell(v_{0}))$ be the maximal interval of existence. Integrating \eqref{eqn quotient} from $c_{m}$ to $\ell(v_{0})$, we obtain  
 \begin{equation}\label{eqn dfn of ell}
    \ell(v_{0})=c_{m}+ \gamma \left(\frac{p-1}{p}\right)^{1/p}\int_{v_{0}}^{\infty}\frac{ds}{\{F(s)-F(v_{0})\}^{1/p}}.
\end{equation}
The quantity $\ell(v_{0})<+\infty$ due to \ref{KO}.
Similarly, let $\Tilde{v}$ be the unique solution of
\begin{equation}\label{ODE2}
   \begin{cases}
     \gamma^{p}\left(|\Tilde{v}'|^{p-1}\Tilde{v}'\right)'=f(\Tilde{v})\quad \text{for all } x\leq c_{m}\\
     \Tilde{v}'(c_{m})=0 \quad \text{and} \quad \Tilde{v}(c_{m})=v_{0}.
 \end{cases} 
\end{equation}
It is given by 
\begin{equation}\label{eqn u2}
    \gamma \left(\frac{p-1}{p}\right)^{1/p}\int_{v_{0}}^{\Tilde{v}(x)}\frac{ds}{\{F(s)-F(v_{0})\}^{1/p}}=x-c_{m}\quad \text{on } (\Tilde{\ell}(v_{0}),c_{m}).
\end{equation}
Define the $C^{1}(\tilde{\ell}(v_{0}),\ell(v_{0}))$ function $w$, where $w(x):=v(x)$ for $x\geq c_{m}$ and $w:=\tilde{v}(x)$ for $x\leq c_{m}$. Then $w$ is the unique solution of \eqref{eqn in one dim without boundary} when $a=\Tilde{\ell}(v_{0})$ and $b =\ell(v_{0})$. By the observation from the beginning of this section, $w$ is convex. In particular, $v_{0}$ is the minimum of $w$.
By uniqueness of the solution $w(x)=w(2c_{m}-x)$, which implies $c_{m}=\frac{\ell(v_{0})+\Tilde{\ell}(v_{0})}{2}$. Thus, up to translations, it is enough to assume that $c_{m}=0$, and consider the interval $(-\ell(v_{0}),\ell(v_{0}))$.

  \textit{Step 2:} The mapping $\ell:(0,\infty)\to(0,\infty)$ is bijective.

Clearly, $\ell$ is injective as the integrand in \eqref{eqn dfn of ell} is positive. Applying change of variables to \eqref{eqn dfn of ell}, we obtain 
        \begin{equation*}
            \ell(t)=\gamma \left(\frac{p-1}{p}\right)^{1/p}\int_{0}^{\infty}\frac{ds}{\left\{F(s+t)-F(t)\right\}^{1/p}}.
        \end{equation*}
        First, $\ell$ is continuous. To show that $\ell$ is onto, it is sufficient to show that $\ell(t)\to \infty$ as $t\to 0$ and $\ell(t)\to 0$ as $t\to \infty$.

        Since $F(s+t)-F(t)\to F(s)$ as $t\to 0$, the Osgood condition \ref{i1} implies that $\ell(t)\to \infty$ as $t\to 0$. 
        
        For the second limit, as $f$ is increasing, we can write
        \begin{equation*}
           F(s+t)-F(t)=\int_{t}^{s+t}f(r)\ dr\geq sf(t), 
        \end{equation*}
    
       Thus 
       \begin{equation*}
           \left\{F(s+t)-F(t)\right\}^{1/p}\geq (sf(t))^{1/p}\to +\infty \text{ as } t\to +\infty.
       \end{equation*}
      Indeed, if $f\leq M$ for some $M\in (0,+\infty)$, we would have $F(t)^{-1/p}\geq (Mt)^{-1/p}$, and so \ref{KO} fails. Thus, $\ell(t)\to 0$ as $t\to \infty$.

      
  \textit{Step 3:}
        Let $c=\frac{a+b}{2}$ and $\delta=\frac{b-a}{2}$. By the previous step, there exists $t_{0}\in (0,\infty)$ such that $\ell(t_{0})=\delta$. The function $\phi$ defined implicitly by 
        \begin{equation*}
              \gamma\left(\frac{p-1}{p}\right)^{1/p}\int_{t_{0}}^{\phi(x)}\frac{ds}{\{F(s)-F(v_{0})\}^{1/p}}=x, \quad \text{for all } x\in (-\delta,\delta),
        \end{equation*}
         solves \eqref{eqn in one dim} for $a=-\delta$ and $b=\delta$. Thus, $v(x)=\phi(x-c)$ is the desired solution.
\end{proof}

 With the assumptions present in the previous theorem, a natural to ask if  $\ell(v_{0})=+\infty$, or equivalently, $v_{0}=0$. The answer is negative. In this regard, the next two results provide the non-existence of a large solution on $\R$.
\begin{lemma}
    Assume \ref{i1}. Let $u$ be the solution of \eqref{eqn in one dim}. Then $u_{0}:=\min_{x}u(x)> 0$.
\end{lemma}
\begin{proof}
    As $u\geq 0$, $u_{0}\geq 0$. If $u_{0}=0$, let $t_{0}:=\min\{t~|~w(t)>0\}$. Then, $t_{0}\geq c_{m}$. For any $t_{0}<x<y$, integrating \eqref{eqn quotient} between $x$ and $y$ implies
    \begin{equation*}
        \gamma \left(\frac{p-1}{p}\right)^{1/p}\int_{w(x)}^{w(y)}\frac{ds}{\left\{(F(s)\right\}^{1/p}}=y-x.
    \end{equation*}
    As $x\to t_{0}^{+}$, the monotone convergence theorem and \ref{i1} gives a contradiction. 
    \end{proof}
    \begin{corollary}
        Let $f$ satisfy \ref{KO} and \ref{i1}. There exist no function $u$ solving \eqref{eqn in one dim} with $a=-\infty$ and $b=\infty$. 
    \end{corollary}
    \begin{proof}
     If such a solution $u$ exists, then by \eqref{eqn dfn of ell} and \ref{KO} $u_{0}(:=\displaystyle\min_{\R}u(x))=0$, a contradiction to the above lemma.
    \end{proof}
    


  \begin{remark}
    In the proof of step 2 of Theorem \ref{thrm existence in one dim with osgood}, the Osgood condition is not used to prove $\ell(t)\to 0$ as $t\to \infty$.
 \end{remark}
 Now we come to the second possibility \ref{i2}. Define 
 \begin{equation*}
     L=\gamma \left(\frac{p-1}{p}\right)^{1/p}\int_{0}^{\infty} \frac{ds}{\left\{F(s)\right\}^{1/p}}
 \end{equation*}
 In view of the proofs of Theorem \ref{thrm existence in one dim with osgood}, one has
 \begin{theorem}\label{thrm existence without osgood}
     Assume \ref{KO}, \ref{a1}, and \ref{i2}. For any $a<b$ in $\R$ with $b-a\leq 2L$, there exists a unique $v$ solving \eqref{eqn in one dim}.
 \end{theorem} 
 \begin{proof}
     The proof follows the same approach as Theorem \ref{thrm existence in one dim with osgood}. the only change is that when $t\to 0$, $\ell(t)\to L$ by  the dominated convergence theorem.
 \end{proof}
 If $b-a>2L$, a solution can be constructed. To this end, up to a translation, we assume that $a=-b$. Let $v_{L}$
be the solution of \eqref{eqn in one dim} for $-a=b=L$ given by Theorem \ref{thrm existence without osgood}. Define the unique solution in this case as follows,
\begin{align*}
      \Tilde{v}(x)=0 \quad \forall~x\in [L-b,b+L],
  \end{align*}
 And it is implicitly defined by
 \begin{align*}
     \gamma\left(\frac{p-1}{p}\right)^{1/p}\int_{0}^{\Tilde{v}(x)}\frac{ds}{\left\{F(s)\right\}^{1/p}}=x-b+L \quad \forall x\in (-b, L-b)\cup (b-L, b).
 \end{align*}
\begin{theorem}
      Assume \ref{KO}, \ref{a1}, and \ref{i2}. For any $a<b$ in $\R$ with $b-a> 2L$, there exists a unique $v$ solving \eqref{eqn in one dim}. Also, $v=0$ in $[L-b+c,b+L+c]$
\end{theorem}
\begin{proof}
    Aided by the above discussion, the solution is given by $v(x)=\Tilde{v}(x-c)$.
\end{proof}
\subsection{Boundary asymptotics:}

\begin{proposition}
    Let $u$ be a solution of \eqref{eqn in one dim}. Then
    \begin{equation*}
       \lim_{x\to b}\frac{\Psi_{H,p}(u(x))}{b-x}= \frac{1}{\gamma},
    \end{equation*}
 and 
  \begin{equation*}
       \lim_{x\to a}\frac{\Psi_{H,p}(u(x))}{x-a}= \frac{1}{\gamma}.
    \end{equation*}
\end{proposition}
\begin{proof}
    Clearly 
    \begin{equation*}
        \left\{F(s)\right\}^{1/p}\geq \left\{F(s)-F(v_{0})\right\}^{1/p}.
    \end{equation*}
   On the other hand, since $F(s)\to \infty$ as $s\to\infty$, given $\epsilon>0$, there exists $s(\epsilon)>0$ such that $F(s)\geq \epsilon F(v_{0})$ for all $s\geq s(\epsilon)$;
   which implies
   \begin{equation*}
       \left\{F(s)-F(v_{0})\right\}^{1/p}\geq \left(1-\frac{1}{\epsilon}\right)^{1/p}\left\{F(s)\right\}^{1/p}
   \end{equation*}
   With these, one can write
   \begin{equation}\label{eqn double inequality}
       \left(1-\frac{1}{\epsilon}\right)^{1/p}\frac{1}{ \left\{F(s)-F(v_{0})\right\}^{1/p}}\leq \frac{1}{\left\{F(s)\right\}^{1/p}}\leq \frac{1}{ \left\{F(s)-F(v_{0})\right\}^{1/p}}.
   \end{equation}
   Since $u(t)\to\infty$ as $t\to b$, we obtain $r>0$ such that for all $t>b-r$, $u(t)\geq s(\epsilon)$. For $t>b-r$ integrating the above inequality from $u(t)$ to $\infty$, we obtain 
   \begin{equation}\label{eqn double integral ineq}
       \gamma\left(\frac{p-1}{p}\right)^{1/p}\left(1-\frac{1}{\epsilon}\right)^{1/p}\int_{u(t)}^{\infty}\frac{1}{ \left\{F(s)-F(v_{0})\right\}^{1/p}}\leq \gamma\Psi_{H,p}(u(t)) \leq \gamma\left(\frac{p-1}{p}\right)^{1/p}\int_{u(t)}^{\infty}\frac{1}{ \left\{F(s)-F(v_{0})\right\}^{1/p}}.
   \end{equation}
   Integrating \eqref{eqn quotient} from $t$ to $b$,  employing  the change of variables $s=v(x)$, and using this in the above inequality we obtain 
   \begin{equation*}
       \left(1-\frac{1}{\epsilon}\right)^{1/p}\leq \frac{\gamma\Psi_{H,p}(u(t))}{b-t}\leq 1.
   \end{equation*}
   As $t\to b$, therefore we obtain  
   \begin{equation*}
       \left(1-\frac{1}{\epsilon}\right)^{1/p}\leq \liminf_{t\to b}\frac{\gamma\Psi_{H,p}(u(t))}{b-t}\leq \limsup_{t\to b}\frac{\gamma\Psi_{H,p}(u(t))}{b-t}\leq 1.
   \end{equation*}
   and also $\epsilon\to \infty$, we obtain 
    \begin{equation*}
     \lim_{t\to b}\frac{\Psi_{H,p}(u(t))}{b-t}= \frac{1}{\gamma}.
   \end{equation*}
   %Also, note that as $\epsilon\to\infty$, $t\to b$. 
   
   Next, as $u(t)\to \infty$ as $t\to a$, there exists $r>0$ such that for all $t<a+r$ we have $u(t)>s(\epsilon)$. The counterpart of \eqref{eqn quotient} for $u=\Tilde{v}$, which solves \eqref{ODE2} is 
   \begin{equation*}
       \gamma\left(\frac{p-1}{p}\right)^{1/p}\frac{-u'(x)}{\left\{F\left(u(x)\right)-F(v_{0})\right\}^{1/p}}=1.
   \end{equation*}
   Integrating the above equation from $a$ to $t$ and using a change of variables $s=v(x)$, we obtain  
   \begin{equation*}
       t-a=\gamma\left(\frac{p-1}{p}\right)^{1/p}\int_{u(t)}^{\infty}\frac{ds}{\left\{F(s)-F(v_{0})\right\}^{1/p}}.
   \end{equation*}
   Using this equation in \eqref{eqn double integral ineq}, we have 
    \begin{equation*}
       \left(1-\frac{1}{\epsilon}\right)^{1/p}\leq \frac{\gamma\Psi_{H,p}(u(t))}{t-a}\leq 1.
   \end{equation*}
   As $\epsilon\to \infty$ we obtain 
   \begin{equation*}
       \lim_{t\to a}\frac{\Psi_{H,p}(u(t))}{t-a}=\frac{1}{\gamma}.\qedhere
   \end{equation*}
\end{proof}
Since $\delta(x)=min\{b-x,x-a\}$ we have the following boundary asymptotics
\begin{theorem}\label{theorem boy asylum 1d}
        Let $u$ be a solution of \eqref{eqn in one dim}. Then
    \begin{equation*}
       \lim_{\delta(x)\to 0}\frac{\Psi_{H,p}(u(x))}{\delta(x)}= \frac{1}{\gamma}.
    \end{equation*}
\end{theorem}

    If we consider the integration with respect to the measure corresponding to $H_{0}$, we can have  
   \begin{proof}[\textbf{Proof of Theorem \ref{dual asym} for $n=1$}]
      Since $H_{0}(x)=\frac{|x|}{\gamma}$, we have $\delta(x)=\gamma \delta_{H_{0}}(x)$. Thus, Theorem \ref{theorem boy asylum 1d} proves Theorem \ref{dual asym}.
   \end{proof}


\section{Existence in higher dimension}\label{higher dim}
 In this section, we discuss the existence of a solution to \eqref{finsler large eqn} in higher dimensions as stated in Theorem \ref{main theorem existence}.  We begin with the proof of the comparison principle stated in Theorem \ref{comparison principle}. 

    \begin{proof}[\textbf{Proof of Theorem \ref{comparison principle}}]
    Suppose not. Then, by continuity, there exists an open ball $B\subset\Omega$ such that $u>v$ on $\overline{B}$. Let $\epsilon_{0}>0$ be such that for every $\epsilon\in (0,\epsilon_{0})$, $0<u-(1+\epsilon)v$ in $B$. Also, by \eqref{boundary condition for comp princ}, let $\delta>0$ be such that $u-(1+\epsilon)v\leq 0$ in $\Omega_{\delta}:=\{x\in\Omega~|~ d(x,\partial\Omega)<\delta\}$ for all $\epsilon\in(0,\epsilon_{0})$. 

On the one hand, since $f$ is strictly increasing, we have
\begin{equation}\label{controduction ineq1}
    \int_{B}\left(f(u)-f(v)\right)(u-v)>0.
\end{equation}
On the other hand, by \eqref{ineq for comp princ}, we have for all $\phi\in W_{0}^{1,p}(\Omega)$
\begin{equation}\label{eqn weak comp ineq}
     \int_{\Omega}\left<\Fin{u}-\Fin{v}, \nabla \phi\right>+(f(u)-f(v))\phi \ dx\leq 0
\end{equation}
Let $\Omega'\subset \Omega\setminus\Omega_{\delta}$ be a subset such that $u-(1+\epsilon)v>0$ in $\Omega'$ and vanishes on $\partial\Omega'$. Then $(u-(1+\epsilon)v)_{+}\in W_{0}^{1,p}(\Omega')$. Define
\begin{equation*}
       I_{\epsilon}:=\int_{B}\left(f(u)-f(v)\right)(u-(1+\epsilon)v)_{+}\ dx.
\end{equation*}
By the definition of $\Omega'$, we have
\begin{equation*}
  I_{\epsilon}  \leq \int_{\Omega'} \left(f(u)-f(v)\right)(u-(1+\epsilon)v)\ dx.
\end{equation*}
Using Lemma \ref{lemma monotone operator} with $x=\nabla u$ and $y=(1+\epsilon)\nabla v$, we can write
\begin{equation*}
    \begin{split}
        I_{\epsilon}\leq  &\int_{\Omega'}\left<\Fin{u}-\Fin{(1+\epsilon)v}, \nabla (u-(1+\epsilon)v)\right>\ dx +\\
        &\quad +\int_{\Omega'} \left(f(u)-f(v)\right)(u-(1+\epsilon)v)\ dx.
    \end{split}
\end{equation*}

Using the homogeneity of $H$, first order binomial expansion on $(1+\epsilon)^{p-1}$ for $p\geq 2$ in the preceding equation. Also, by the inequality \eqref{eqn weak comp ineq} and $(4)$ of Lemma \ref{lemma_properties of H}, one has
\begin{equation*}
    \begin{split}
        I_{\epsilon}
        \leq & \int_{\Omega'}\left<\Fin{u}-(1+(p-1)\epsilon+o(\epsilon))\Fin{v}, \nabla (u-(1+\epsilon)v)\right>\ dx\\
        &+\int_{\Omega'} \left(f(u)-f(v)\right)(u-(1+\epsilon)v)\ dx\\
        \leq &-\left((p-1)\epsilon+o(\epsilon)\right) \int_{\Omega'}\left<\Fin{v}, \nabla(u-(1+\epsilon)v)\right>\ dx\\
        \leq & -\left((p-1)\epsilon +o(\epsilon)\right)\left\{\int_{\Omega'}\left<\Fin{v}, \nabla u \right>\ dx -(1+\epsilon) \int_{\Omega'}H^{p}(\nabla v)\ dx\right\} .
    \end{split}
\end{equation*}
As $\epsilon\to 0$, applying Fatou's lemma to the definition of $I_{\epsilon}$ and then using the above inequality, we obtain 
\begin{equation}\label{contraduction ineq2}
    \int_{B}\left(f(u)-f(v)\right)(u-v)\ dx\leq \liminf_{\epsilon\to 0}I_{\epsilon}\leq 0.
\end{equation}
A contradiction to \eqref{controduction ineq1}.  Thus, $u\leq v$ on $\Omega$. 
  \end{proof}

Since $\Omega$ is an open set, given $x\in \Omega$ choose $R>0$ such that $B_{R}(x)\subset \Omega$. Since $|x-y|\leq \theta_{2}H_{0}(x-y)$, take any $r\leq R/\theta_{2}$, then $\mathcal{W}_{r}(x)\subset B_{R}(x)\subset\Omega$. Using this, we construct a local barrier to a large solution, which is radial on a Wulff ball.  
  \begin{proposition}\label{prop_weak_local_bound}
    Let $\Omega$ be a bounded domain in $\rn$, $n\geq 2$. Assume \ref{KO}, \ref{a1}. Let $u\in W_{loc}^{1,p}(\Omega)\cap C(\Omega)$ be a weak subsolution of \eqref{finsler large eqn}. For any $x_{0}\in \Omega$ and $R>0$ such that $\mathcal{W}_{R}(x_{0})\subset \subset \Omega$, we have
    \begin{equation}\label{eqn_local_bound}
        u(x)\leq \omega\Big(\frac{R}{2}\Big) \quad \text{for all } x\in \mathcal{W}_{\frac{R}{2}}(x_{0}).
    \end{equation}
    Where, $\omega$ is a solution of
     \begin{equation}\label{eqn omega is a large solution}
        \begin{cases}
            \left(|\omega'(t)|^{p-2}\omega'(t)\right)'=f(\omega(t)) \quad &\text{in } (0,2R)\\
            \omega(t)\to\infty \quad &\text{as } t\to 0, 2R.
        \end{cases}
    \end{equation}
\end{proposition}
 \begin{proof}
     Using $\omega$, we construct a radially symmetric weak boundary blow-up  supersolution $v$ of \eqref{finsler eqn} in $\mathcal{W}_{R}(x_{0})$. Then the bound is obtained by Theorem \ref{comparison principle} and the radial symmetry of $v$.

     
    Define $v(x)=\omega(R-H_{0}(x-x_{0}))\in C(\mathcal{W}_{R}(x_{0}))$. For any $\phi\in C_{c}^{\infty}(\mathcal{W}_{R}(x_{0}))$, $\phi\geq 0$. Setting $x=x_{0}+\frac{R-t}{R}\theta(x)$, where $H_{0}(x-x_{0})=R-t$ and $\theta(x)\in \partial \mathcal{W}_{R}(x_{0})$, we obtain  by Lemma \ref{lemma_properties of H}, Lemma \ref{H0} and a change of variables ($ R-H_{0}(x-x_{0})\to t$) that
     \begin{equation}\label{eqn_eta'}
     \begin{split}
    &\int_{\mathcal{W}_{R}(x_{0})}\left<\Fin{v}, \nabla\phi \right>\ dx\\
    &=-\int_{\mathcal{W}_{R}(x_{0})}\left<|\omega'(R-H_{0}(x-x_{0})|^{p-2}\omega'(R-H_{0}(x-x_{0})) \frac{x-x_{0}}{H_{0}(x-x_{0})}, \nabla \phi \right>\ dx\\
      &=\int_{0}^{R}\int_{\partial \mathcal{W}_{R-t}(x_{0})}\left<|\omega'(t)|^{p-2}\omega'(t) \frac{\theta(x)}{R}, \nabla \phi(x) \right> \ \frac{d\mathcal{H}^{n-1}(x)}{|\nabla H_{0}(x-x_{0})|} dt\\
        &=\int_{0}^{R}|\omega'(t)|^{p-2}\omega'(t)~ \int_{\partial \mathcal{W}_{R}(x_{0})}-\frac{\partial \phi}{\partial t}\left(x_{0}+\frac{R-t}{R}\theta(x)\right) \frac{(R-t)^{n-1}}{R^{n-1}} \ \frac{d\mathcal{H}^{n-1}(x)}{|\nabla H_{0}(\theta(x))|} dt.\\ 
         \end{split}
     \end{equation}
    Where we have by the chain rule that $-\frac{\partial \phi}{\partial t}(t,x)=\frac{\theta(x)}{R}\cdot \nabla \phi(t,x)$. %and $\nabla H_{0}(sx)=\nabla H_{0}(x)$ for all $s\geq 0$.\\
     
    Notice that, as $\phi\geq 0$ and $s<R$, 
    \begin{equation*}
    \begin{split}
        \frac{\partial \phi}{\partial s}(s,x) (R-s)^{n-1}&=\frac{\partial}{\partial s}\left(\phi(s,x)(R-s)^{n-1}\right)+(n-1)\phi(s,x)(R-s)^{n-2}\\
        &\geq \frac{\partial}{\partial s}\left(\phi(s,x)(R-s)^{n-1}\right).
        \end{split}
    \end{equation*}
    And, $\omega$ is a solution of equation \eqref{eqn omega is a large solution} implies $\omega'\leq 0$ on $(0,R)$. Therefore, by \eqref{eqn_eta'}, we obtain  
    \begin{equation*}
        \begin{split}
            &\int_{\mathcal{W}_{R}(x_{0})}\left<\Fin{v}, \nabla\phi \right>\ dx\geq \\
            &\geq \int_{0}^{R}|\omega'(t)|^{p-2}\omega'(t)~ \int_{\partial \mathcal{W}_{R}(x_{0})}-\frac{\partial }{\partial t}\left(\phi(t,x) \frac{(R-t)^{n-1}}{R^{n-1}}\right) \ \frac{d\mathcal{H}^{n-1}(x)}{|\nabla H_{0}(\theta(x))|} dt.\\
            &\geq -\int_{\partial \mathcal{W}_{R}(x_{0})}\int_{0}^{R}f(\omega(t))\phi\left(x_{0}+\frac{R-t}{R}\theta(x)\right)\frac{(R-t)^{n-1}}{R^{n-1}}\ dt \ \frac{d\mathcal{H}^{n-1}(x)}{|\nabla H_{0}(\theta(x))|}\\
            &\geq -\int_{0}^{R} \int_{\partial \mathcal{W}_{R-t}(x_{0})}f(\omega(t))\phi(t,x)\ \frac{d\mathcal{H}^{n-1}(x)}{|\nabla H_{0}(\theta(x))|}\ dt\\
            &\geq -\int_{\mathcal{W}_{R}(x_{0})}f(v(x))\phi(x)\ dx.
        \end{split}
    \end{equation*}
    Also,$v(x)\to \infty$ as $H_{0}(x-x_{0})\to R$.
Since $u<\infty$ in $\mathcal{W}_{R}(x_{0})$, Theorem \ref{comparison principle} implies 
    \begin{equation*}
        u(x)\leq v(x)\quad \text{for all }x\in \mathcal{W}_{R}(x_{0}).
    \end{equation*}
    This and the monotonicity of $\omega$ in $(0,R)$ leads to \eqref{eqn_local_bound}. 
 \end{proof}
 \begin{remark}
     The argument in the above proposition is independent of the boundary values. Thus, the same holds for solutions with finite Dirichlet boundary values. 
 \end{remark}
 The next proposition gives the local $L^{p}$ bound for the gradient of the solution.
 \begin{proposition}\label{prop_grad lp bound}
     Let $\Omega$ be a bounded domain and let $u\in W_{loc}^{1,p}(\Omega)\cap C(\Omega)$ is the solution of \eqref{finsler large eqn} then there exists a positive constant $C=C(p,K, \theta_{1},\theta_{2})$ for any compact set $K$ in $\Omega$ such that  
     \begin{equation*}
         \|\nabla u\|^{p}_{L^{p}(K)}\leq C \left(1+f\left(\|u\|_{L^{\infty}(\Tilde{K})}\right)\right)\|u\|_{L^{\infty}(\Tilde{K})}.
     \end{equation*}
     for some compact set $\Tilde{K}$ such that $K\subset \Tilde{K}\subset \Omega$.
 \end{proposition}
 \begin{proof}
     Let $\phi\in C_{c}^{\infty}(\Tilde{K})$, $0\leq\phi\leq 1$ and $\phi=1$ on $K$. Substituting $u\phi^{p}$ in the weak formulation and using Lemma \ref{lemma_properties of H}, we obtain
     \begin{equation*}
         \begin{split}
             \theta_{1}^{p}\int_{\Tilde{K}}|\nabla u|^{p}\phi^{p} &\leq \int_{\Tilde{K}}H^{p}(\nabla u) \phi^{p}\ dx=\int_{\Tilde{K}}f(u)u\phi^{p}\ dx-p\int_{\Tilde{K}}u\phi^{p-1}\left<\Fin{u}, \nabla \phi\right>\ dx\\
             &\leq f(\|u\|_{L^{\infty}(\Tilde{K})})\|u\|_{L^{\infty}(\Tilde{K})}|\Tilde{K}|+p\theta_{2}^{p-1}\|u\|_{L^{\infty}(\Tilde{K})} \int_{\Tilde{K}}\phi^{p-1}  |\nabla u|^{p-1}|\nabla\phi|\ dx.
         \end{split}
     \end{equation*}
     Now let $\epsilon>0$. Applying Young's inequality with the parameter $\epsilon$ for the second integral on the right side
     \begin{equation*}
         \int_{\Tilde{K}}\phi^{p-1}  |\nabla u|^{p-1}|\nabla\phi|\ dx\leq \frac{\epsilon(p-1)}{p}\int_{\Tilde{K}}\phi^{p}|\nabla u|^{p}\ dx +\frac{1}{p\epsilon^{p-1}}\int_{\Tilde{K}}|\nabla\phi|^{p}\ dx.
     \end{equation*}
     Choosing $\epsilon=\frac{\theta_{1}^{p}p}{2(p-1)}$, we obtain 
     \begin{equation*}
         \int_{K}|\nabla u|^{p}\ dx\leq C(p,K,\theta_{1},\theta_{2})\|u\|_{L^{\infty}(\Tilde{K})}\left(f(\|u\|_{L^{\infty}(\Tilde{K})})+1\right).\qedhere
     \end{equation*}
 \end{proof}

 Next, we address the existence:
 \begin{proof}[\textbf{Proof of Theorem \ref{main theorem existence} for $n\geq 2$}]
   The idea is to get the large solution as a limit of a sequence of solutions with finite boundary data.
   
  \textit{Step 1:} Let $u_{k}$ be the solution to \eqref{dirchlet problem} as in Definition \ref{dfn weak solution to dirchlet prob} with $g=k$. $\{u_{k}\}_{k}$ is an increasing sequence by Theorem \ref{comparison principle} and the sequence is locally uniformly bounded by Proposition \ref{prop_weak_local_bound}. For any $x\in\Omega$, define
   \begin{equation}\label{dfn of u}
       u(x):=\lim_{k\to\infty}u_{k}(x).
   \end{equation}
   By Proposition \ref{prop_weak_local_bound}, $u\in L_{loc}^{\infty}(\Omega)$. Then, by the dominated convergence theorem $u_{k}\to u$ in $L^{q}_{loc}(\Omega)$ for all $1\leq q<\infty$. Further, by Proposition \ref{prop_weak_local_bound}
 and Proposition \ref{prop_grad lp bound}, the sequence $\{u_{k}\}_{k}$ is bounded in $W^{1,p}_{loc}(\Omega)$ and hence by \eqref{dfn of u} and Sobolev embedding, $u_{n_{k}}\rightharpoonup u$ in $W_{loc}^{1,p}(\Omega)$. Thus $u\in W_{loc}^{1,p}(\Omega)$. 
 
 Next, we show that this candidate $u$ is the required solution.
 
\textit{Step 2:}
 The continuity of $f$ implies $f(u_{k}(x))\to f(u(x))$ for all $x\in\Omega$. The monotonicity of $f$ and the monotone convergence theorem imply $f(u_{k})\to f(u)$ in $L^{1}_{loc}(\Omega)$, thus $f(u)\in L^{1}_{loc}(\Omega)\cap L^{\infty}_{loc}(\Omega)$.  

 Let $\Omega'\subset \subset \Omega$. For any $\phi\in W_{0}^{1,p}(\Omega')$, by the dominated convergence theorem 
 \begin{equation*}
     \int_{\Omega'}f(u_{k})\phi\ dx\to \int_{\Omega'}f(u)\phi\ dx.
 \end{equation*}

\textit{Step 3:}
     For any $v\in W^{1,p}(\Omega')$, by Lemma \ref{lemma_properties of H} and the H\"older inequality, we obtain 
    \begin{equation*}
    \begin{split}
         \int_{\Omega'}\left<\Fin{u_{n_{k}}}, \nabla v\right> \ dx & \leq C \int_{\Omega'}H^{p-1}(\nabla u_{n_{k}}) |\nabla v| \ dx \\
         & \leq C\theta_{2}^{p-1}\int_{\Omega'} |\nabla u_{n_{k}}|^{p-1}|\nabla v|\ dx \\
         &\leq  C\theta_{2}^{p-1} \|\nabla u_{n_{k}}\|^{p-1}_{L^{p}(\Omega')} \|\nabla v\|_{L^{p}(\Omega')}.
    \end{split}
    \end{equation*}
    Therefore, by Propositions \ref{prop_grad lp bound} and Propositions \ref{prop_weak_local_bound}, the sequence $\mathcal{T}_{k}\in \left(W^{1,p}(\Omega')\right)'$ given by 
    \begin{equation*}
        \left<\mathcal{T}_{k}, v\right>:=\int_{\Omega'}\left<\Fin{u_{n_{k}}} , \nabla v\right>\ dx
    \end{equation*}
    is bounded and hence there exists a subsequence denoted by $u_{n_{k}}$ and a $\mathcal{T}\in \left(W^{1,p}(\Omega')\right)'$ such that 
    \begin{equation}\label{eqn weak convergence in dual}
        \int_{\Omega'} \left<\Fin{u_{n_{k}}}, \nabla v\right> \ dx=\left<\mathcal{T}_{n_{k}}, v\right> \to \left<\mathcal{T}, v\right> \quad \text{for all } v\in  W^{1,p}(\Omega').
    \end{equation}
    Moreover, for $\phi\in W_{0}^{1,p}(\Omega')$
    \begin{equation*}
        \int_{\Omega'} \left<\Fin{u_{n_{k}}}, \nabla \phi \right>\ dx= -\int_{\Omega'}f(u_{n_{k}})\phi \ dx \to -\int_{\Omega'}f(u)\phi \ dx,
    \end{equation*}
    as $n_{k}\to \infty$. Thus, 
    \begin{equation}\label{L equals to f}
        \left<\mathcal{T}, \phi\right>=-\int_{\Omega'}f(u)\phi \ dx \quad \text{for all } \phi\in W_{0}^{1,p}(\Omega').
    \end{equation}

    \textit{Step 4:}
    Finally, we will show that
    \begin{equation*}
        \left<\mathcal{T}, \phi\right>=\int_{\Omega'} \left<\Fin{u}, \nabla \phi \right> \ dx \quad \text{for all } \phi \in W_{0}^{1,p}(\Omega'),
    \end{equation*}
     using Browder-Minty type argument (cf. \cite{Zeidler1990,evans}).
    
    By (monotonicity) Lemma \ref{lemma monotone operator}
    \begin{equation*}
        \int_{\Omega'} \left<\Fin{u_{n_{k}}}-\Fin{v}, \nabla(u_{n_{k}}-v) \right> \ dx \geq 0, ~~ \text{for any } v \in W^{1,p}(\Omega').
   \end{equation*}
   Further, as $n_{k}\to \infty$, by $u_{n_{k}}\rightharpoonup u$, $\mathcal{T}_{n_{k}}\overset{*}\rightharpoonup \mathcal{T}$, \eqref{eqn weak convergence in dual} and a diagonal argument we obtain    
   \begin{equation*}
       \left<\mathcal{T}-\int_{\Omega'}\Fin{v}, u-v\right>\geq 0.
   \end{equation*}
   Now, for any $w\in W^{1,p}(\Omega')$ substitute $v=u-tw$, $t\in [0,1]$ to get 
   \begin{equation}\label{BM}
       \left<\mathcal{T}-\int_{\Omega'}\Fin{u-tw}, w\right>\geq 0.
   \end{equation}
   Also, by using $H\in C^{2}(\rn\setminus\{0\})$ when $\nabla u\neq 0$ and the 1-homogeneity of $H$ otherwise; the map $t\longmapsto \int_{\Omega'}\left<\Fin{u-tw}, \nabla w\right>$ is  continuous. For any $w\in W^{1,p}(\Omega')$, $t\to 0$ in \eqref{BM} implies 
   \begin{equation*}
       \left<\mathcal{T}-\int_{\Omega'}\Fin{u}, w\right>\geq 0.
   \end{equation*}
   Replacing $w$ with $-w$ in the above inequality together with \eqref{L equals to f} , we obtain 
   \begin{equation}
       \int_{\Omega'}\left<\Fin{u}, \nabla \phi \right>\ dx = -\int_{\Omega'}f(u)\phi \ dx \quad \text{for all } \phi\in W_{0}^{1,p}(\Omega').
   \end{equation}
   This establishes the result.
\end{proof}

\section{Boundary Asymptotics}\label{asym}
In this section, we prove Theorem \ref{dual asym} in two steps. One is Lemma \ref{lemma asym 1} and the other is Lemma \ref{lemma asym 2}. In order to prove Lemma \ref{lemma asym 2}, we are led to prove the existence as well as the boundary behaviour of a radially symmetric (with respect to the dual norm ${H_{0}}$) solution on an annulus defined with respect to $H_{0}$. Since Theorem \ref{dual asym} is proved for $n=1$ in Section \ref{1d}, we assume $n\geq 2$ throughout this section.

  Let $\Omega\subset\R^{n}$, be a bounded domain with $C^{2}$ boundary. Let $\delta_{H_{0}}(x)=\inf_{z\in \partial\Omega}H_{0}(x-z)$ denote the Minkowski distance from the boundary of $\Omega$. We first collect some of the properties of $\delta_{H_{0}}$ found in \cite{CrastaMalusa2007}. 

Since $\Omega$ is $C^{2}$, there exists $\mu>0$ such that $\delta_{H_{0}}$ is $C^{2}$ on $\Omega_{\mu}:=\{x\in\Omega~|~\delta_{H_{0}}(x)<\mu\}$, and moreover, for any $x\in \Omega_{\mu}$, there is a unique $z(x)\in \partial\Omega$ such that
 \begin{equation}\label{rep of x in Omega mu}
 \begin{split}
    &\delta_{H_{0}}(x)=H_{0}(x-z(x)) \quad\text{and}\\
   & x=z(x)+\delta_{H_{0}}(x)\nabla H(\nabla\delta_{H_{0}}(x)).
 \end{split}
       \end{equation} 
       %Also, from Lemma \ref{H0}, $H_{0}\left(\nabla H(\nabla\delta_{H_{0}}(x))\right)=1$.
Further, there exists $R>0$, such that given any $z\in \partial\Omega$, there exists $x^{int}(z)\in \Omega$ and $x^{ext}(z)\in \R^{n}\setminus\Omega$, such that $\mathcal{W}_{R}(x^{int}(z))\cap \partial\Omega=\{z\}$ and $\mathcal{W}_{R}\left(x^{ext}(z)\right)\cap \partial\Omega=\{z\}$ (By \cite[Remark 4.2, Proposition 3.3 and also Proposition 4.6]{CrastaMalusa2007}. Moreover,
\begin{equation}\label{rep of center of ball}
    x^{int}=z+R\nabla H\left(\nabla \delta_{H_{0}}(x^{int})\right)\quad \text{and}\quad x^{ext}= z-R\nabla H\left(\nabla \delta_{H_{0}}(x^{ext})\right)
\end{equation}

Owing to Theorem \ref{theorem boy asylum 1d} and the construction of the super solution $v$ in Proposition \ref{prop_weak_local_bound}, we have the following result.
   \begin{lemma}\label{lemma asym 1}
       Let $\Omega$ be a bounded domain in $\rn$ which has the uniform interior ball condition with respect to $H_{0}$. Let $u$ be a solution of \eqref{finsler large eqn}.  Then 
       \begin{equation*}
           \liminf_{\delta_{H_{0}}(x)\to 0}\frac{\Psi_{H,p}\left(u(x)\right)}{\delta_{H_{0}}(x)}\geq 1.
       \end{equation*}
   \end{lemma}
   \begin{proof}
   Let $x\in \Omega_{\mu}$, then there exists $z(x)\in \partial\Omega$ such that \eqref{rep of x in Omega mu} holds. Denoting $z=z(x)$, let $\mathcal{W}_{R}(x^{int})$ be an interior $H_{0}$-ball associated with $z$ and let $\omega_{r}$ be the solution of \eqref{eqn omega is a large solution} on $(0,2(R-r))$ for $0<r<R$.
   
   Let $v_{r}(x)=\omega_{r}\left(R-r-H_{0}(x-x^{int})\right)$ be defined on $\mathcal{W}_{R-r}(x^{int})$. As in the proof of Proposition \ref{prop_weak_local_bound}, $u\leq v_{r}$ on $\mathcal{W}_{R-r}(x^{int})$. By Theorem \ref{comparison principle}, $\{\omega_{r}(x)\}_{r\leq r_{0}}$ is a decreasing sequence for a fixed $x\in \mathcal{W}_{R}(x^{int})$ and $r_{0}$ sufficiently small. Moreover, this sequence is pointwise bounded below by the solution of \eqref{eqn omega is a large solution} on $(0,2R)$. As $\omega_{r_{0}}$ is convex and hence $\omega_{r_{0}}$ is locally bounded, by the dominated convergence theorem, $\omega_{r}\to \omega:=\inf_{r}\omega_{r}$ in $L^{p}_{loc}(0,2R)$ for $r>r_{0}$. By Proposition \ref{prop_grad lp bound}, there exists a subsequence denoted by $\omega_{r}$ such that $\omega_{r}\rightharpoonup \omega$ in $W^{1,p}_{loc}(0,2R)$. As in the proof of Theorem \ref{main theorem existence}, one obtains that $\omega(x)$ is the solution of \eqref{eqn omega is a large solution} on $(0,2R)$.  
    
    Further, since $u\leq \omega(R-H_{0}(x-x^{int}))$ by the above discussion, and $\Psi_{H,p}$ is decreasing, one has $\Psi_{H,p}(u(x))\geq \Psi_{H,p}(\omega(R-H_{0}(x-x^{int})))$ for all $x\in \mathcal{W}_{R}(x^{int})$. By \eqref{rep of x in Omega mu} and \eqref{rep of center of ball}, we obtain
       \begin{equation*}
           \frac{\Psi_{H,p}\left(u(x)\right)}{\delta_{H_{0}}(x)}=\frac{\Psi_{H,p}\left(u(x)\right)}{H_{0}(x-z)}\geq \frac{\Psi_{H,p}\left(\omega\left(R-H_{0}(x-x^{int})\right)\right)}{H_{0}(x-z)}
           = \frac{\Psi_{H,p}\left(\omega\left(R-H_{0}(x-x^{int})\right)\right)}{R-H_{0}(x-x^{int})}.
            \end{equation*}
      Given $\epsilon>0$, by Theorem \ref{dual asym} for $n=1$ (proved in Section \ref{1d}), there exists $\delta>0$ such that for all $x\in \mathcal{W}_{R}(x^{int})$ with $R-H_{0}(x-x^{int})<\delta$,
      \begin{equation*}
        \frac{\Psi_{H,p}(\omega\left(R-H_{0}(x-x^{int}))\right)}{R-H_{0}(x-x^{int})}\geq 1-\epsilon . 
      \end{equation*}
 Thanks to the uniform ball condition. Let $\mu$ be as in \eqref{rep of x in Omega mu}, choosing $\delta\leq \mu$, for all $x\in \Omega_{\delta}:=\{x\in \Omega~|~\delta_{H_{0}}(x)<\delta\}$,  we infer 
\begin{equation*}
    \frac{\Psi_{H,p}\left(u(x)\right)}{\delta_{H_{0}}(x)}\geq 1-\epsilon.
\end{equation*}
      Implies that 
      \begin{equation*}
          \liminf_{\delta_{H_{0}}(x)\to 0} \frac{\Psi_{H,p}\left(u(x)\right)}{\delta_{H_{0}}(x)}\geq 1.\qedhere
      \end{equation*}
   \end{proof}

    To prove Theorem \ref{dual asym}, it remains to establish the reverse inequality involving the limit superior. To this end, we consider the annulus $A_{R_{1}}^{R_{2}}(x_{0})$ defined with respect to the dual norm $H_{0}$, centered at $x_{0}$, with inner radius $R_{1}$ and outer radius $R_{2}$. Take 
   \begin{equation}\label{eqn annulus}
       \begin{cases}
           \Delta_{H}^{p}u=f(u)\quad &\text{in }A_{R_{1}}^{R_{2}}(x_{0})\\
           u(x)\to \infty &\text{as } x\to R_{1}\\
           u(x)\to 0 &\text{as } x\to R_{2}.
       \end{cases}
   \end{equation}
   Next theorem deals with the existence of a radial solution to the above problem.
   \begin{theorem}\label{existence on annulus}
       The equation \eqref{eqn annulus} admits a non-negative radially symmetric solution in $W_{loc}^{1,p}(A_{R_{1}}^{R_{2}}(x_{0}))$.
   \end{theorem}
   \begin{proof}
       Without loss of generality, we assume that the annulus is centred at the origin. Let $v$ be the solution of \eqref{eqn annulus}, and suppose $v(x)=w(H_{0}(x))$ for some function $w:[0,\infty)\to [0,\infty)$. Then for any $\phi\in C_{c}^{\infty}(A_{R_{1}}^{R_{2}}(0))$ such that $\phi(x)=\psi(H_{0}(x))$,
       \begin{equation*}
           -\int_{A_{R_{1}}^{R_{2}}(0)} \left<|w'(H_{0}(x))|^{p-1}H^{p-1}(\nabla H_{0}(x))\frac{w'(H_{0}(x))}{|w'(H_{0}(x))|}\nabla H(\nabla H_{0}(x)),\psi' (H_{0}(x))\nabla H_{0}(x)\right>\ dx=\int_{A_{R_{1}}^{R_{2}}(0)}f(v)\phi \ dx.
       \end{equation*}
       Using Lemma \ref{lemma_properties of H}\ and \ref{H0},
       \begin{equation*}
                      -\int_{A_{R_{1}}^{R_{2}}(0)} |w'(H_{0}(x))|^{p-2} w'(H_{0}(x)) \psi' (H_{0}(x)) \ dx=\int_{A_{R_{1}}^{R_{2}}(0)}f(w(H_{0}(x))\psi(H_{0}(x)) \ dx.
       \end{equation*}
       Changing variables, we obtain 
       \begin{equation*}
           \int_{R_{1}}^{R_{2}}\int_{B_{t}(0)}|w'(t)|^{p-2}w'(t)\psi'(t) \frac{\ d\mathcal{H}(x)}{|\nabla H_{0}(x)|}\ dt=\int_{R_{1}}^{R_{2}}\int_{B_{t}(0)}f(w(t))\psi(t) \frac{\ d\mathcal{H}(x)}{|\nabla H_{0}(x)|}\ dt.
       \end{equation*}
       Which implies
       \begin{equation*}
           \int_{R_{1}}^{R_{2}} t^{n-1}|w'(t)|^{p-2}w'(t)\psi'(t) \ dt=\int_{R_{1}}^{R_{2}} t^{n-1}f(w(t))\psi(t) \ dt.
       \end{equation*}
       In other words, $w$ solves 
       \begin{equation}\label{eqn radial solution}
           \begin{cases}
               \left(t^{n-1}|w'(t)|^{p-2}w'(t)\right)'=t^{n-1}f(w(t)) \quad\text{ in } (R_{1},R_{2})\\
               w(t)\to \infty \text{ as }t\to R_{1} \text{ and }w(t)\to 0\text{ as }t\to R_{2}.
           \end{cases}
       \end{equation}
       To construct a solution to this, let $k\in \mathbb{N}$ and consider the problem
       \begin{equation*}
           \begin{cases}
               \left((R_{2}-t)^{n-1}|\tilde{w}'(t)|^{p-2}\tilde{w}'(t)\right)'=(R_{2}-t)^{n-1}f(\tilde{w}(t)) \quad\text{ in } (0,R_{2}-R_{1})\\
               \tilde{w}(0)=0 \text{ and }\tilde{w}(R_{2}-R_{1})=k,
           \end{cases}
       \end{equation*}
       which has a $C^{1}(0,R_{2}-R_{1})$ solution $\tilde{w}_{k}\geq 0$ by \cite[Proposition 4.2.1]{PucciSerrinbook2007}. Then $v_{k}(x)=w_{k}(H_{0}(x))=\tilde{w}_{k}(R_{2}-H_{0}(x))$ for $x\in A_{R_{1}}^{R_{2}}(0)$ is the solution of
       \begin{equation*}
           \begin{cases}
           \Delta_{H}^{p}v_{k}=f(v_{k})\quad &\text{in }A_{R_{1}}^{R_{2}}(0)\\
           v_{k}(x)\to k &\text{as } x\to R_{1}\\
           v_{k}(x)\to 0 &\text{as } x\to R_{2}.
       \end{cases}
       \end{equation*}
       The sequence $\{v_{k}\}_{k}$ is an increasing by Theorem \ref{comparison principle} and is locally bounded by Proposition \ref{prop_weak_local_bound}. Define $v(x):=\sup_{k}v_{k}(x)$. Then, since Proposition \ref{prop_grad lp bound} also holds for $v_{k}$, as in the proof of Theorem \ref{main theorem existence}, $v$ is the required non-negative radial solution.
   \end{proof}
   \begin{remark}
     Under the assumption \ref{i1} $\tilde{w}_{k}'>0$ by \cite[Proposition 4.2.2]{PucciSerrinbook2007}.  Since $w_{k}(t)=\tilde{w}_{k}(R_{2}-t)$ and $w(H_{0}(x))=v(x)\lim v_{k}(x)=\lim w_{k}(H_{0}(x))$, we obtain that $w$ is decreasing by taking the limit as $k\to \infty$ in $w_{k}(t_{1})\geq w_{k}(t_{2})$ whenever $t_{1}<t_{2}$.
   \end{remark}
   
   Next, we look at the asymptotic of the solution of \eqref{eqn annulus} as $t\to R_{1}$.
   \begin{proposition}\label{prop annuls asym}
       Assume \ref{KO} and \ref{i1}. Let $v\in W_{loc}^{1,p}(A_{R_{1}}^{R_{2}}(0))$ be the radially symmetric solution of \eqref{eqn annulus}. Then 
       \begin{equation*}
           \limsup_{|x|\to R_{1}}\frac{\Psi_{H,p}\left(v(x)\right)}{H_{0}(x)-R_{1}}\leq 1.
       \end{equation*}
   \end{proposition}
\begin{proof}
The idea of the proof is as follows: We consider $\Psi_{H,p}$ with $r=w(t):=v(H_{0}(x))$ and apply a change of variable to obtain $\int_{t}^{\infty}w'(s)/\{F(w(s)\}^{1/p}~ds$. Using the equation and the symmetry, one can compare the quantities $w'$ and $F(w)$ present in the integral, to obtain the desired inequality.  


Let $v(x)=w(H_{0}(x))$. Then $w$ solves \eqref{eqn radial solution}, or equivalently 
\begin{equation*}
    \left(|w'(t)|^{p-2}w'(t)\right)'+\frac{(n-1)}{t}|w'(t)|^{p-2}w'(t)=f(w(t)).
\end{equation*}
Since $w'<0$, the above equation can be written as
\begin{equation*}
    \left(-(-w'(t))^{p-1}\right)'-\frac{(n-1)}{t}(-w'(t))^{p-1}=f(w(t)).
\end{equation*}
Which implies
\begin{equation*}
    (p-1)(-w'(t))^{p-2}w''(t)-\frac{(n-1)}{t}(-w'(t))^{p-1}=f(w(t)).
\end{equation*}
    Multiplying by $\frac{p}{p-1}(-w'(t))t^{\frac{p(n-1)}{p-1}}$, we obtain 
    \begin{equation*}
        \left(t^{\frac{p(n-1)}{p-1}}|w'(t)|^{p}\right)'=\frac{p}{p-1}t^{\frac{p(n-1)}{p-1}}w'(t)f(w(t)).
    \end{equation*}
    Let $R_{1}<t<t_{0}<R_{2}$. Integrating the above equation from $t$ to $t_{0}$,
    \begin{equation*}
        t_{0}^{\frac{p(n-1)}{p-1}}|w'(t_{0})|^{p}-t^{\frac{p(n-1)}{p-1}}|w'(t)|^{p}=\frac{p}{p-1}\int_{t}^{t_{0}}s^{\frac{p(n-1)}{p-1}}dF(w(s)).
    \end{equation*}
    By adding and subtracting $\frac{p}{p-1}t^{\frac{p(n-1)}{p-1}}\int_{t}^{t_{0}}dF(w(s))$,
   % \begin{equation*}
    %    \begin{split}
    %        t^{\frac{p(n-1)}{p-1}}|w'(t)|^{p}=t_{0}^{\frac{p(n-1)}{p-1}}|w'(t_{0})|^{p}-&\frac{p}{p-1}t^{\frac{p(n-1)}{p-1}}\int_{t}^{t_{0}}dF(w(s))+\frac{p}{p-1}t^{\frac{p(n-1)}{p-1}}\int_{t}^{t_{0}}dF(w(s))-\\
    %        &-\frac{p}{p-1}\int_{t}^{t_{0}}s^{\frac{p(n-1)}{p-1}}dF(w(s)).
    %    \end{split}
    %\end{equation*}
    %Implies
    we obtain
    \begin{equation*}
        \begin{split}
            |w'(t)|^{p}=&\left(\frac{t_{0}}{t}\right)^{\frac{p(n-1)}{p-1}}|w'(t_{0})|^{p}-\frac{p}{p-1}\left\{F(w(t_{0}))-F(w(t))\right\}-\frac{p}{p-1}\int_{t}^{t_{0}}\bigg\{\left(\frac{s}{t}\right)^{\frac{p(n-1)}{p-1}}-1\bigg\}dF(w(s))\\
            &\leq \left(\frac{t_{0}}{t}\right)^{\frac{p(n-1)}{p-1}}|w'(t_{0})|^{p}+\frac{p}{p-1}F(w(t))+\frac{p}{p-1}\int_{t_{0}}^{t}\bigg|\left(\frac{s}{t}\right)^{\frac{p(n-1)}{p-1}}-1\bigg|dF(w(s))\\
            &\leq \left(\frac{t_{0}}{R_{1}}\right)^{\frac{p(n-1)}{p-1}}|w'(t_{0})|^{p}+\frac{p}{p-1}F(w(t))+\frac{p}{p-1}\bigg(\left(\frac{t_{0}}{R_{1}}\right)^{\frac{p(n-1)}{p-1}}-1\bigg)\left\{F(w(t))-F(w(t_{0}))\right\}\\
            &\leq \left(\frac{t_{0}}{R_{1}}\right)^{\frac{p(n-1)}{p-1}}|w'(t_{0})|^{p}+\frac{p}{p-1}F(w(t))+\frac{p}{p-1}\bigg(\left(\frac{t_{0}}{R_{1}}\right)^{\frac{p(n-1)}{p-1}}-1\bigg)F(w(t)).
            \end{split}
    \end{equation*}
    Given any $\epsilon>0$, choose $R_{1}<t_{0}<R_{2}$ such that for all $R_{1}<t\leq t_{0}$
    \begin{equation*}
        \left(\frac{t_{0}}{R_{1}}\right)^{\frac{p(n-1)}{p-1}}|w'(t_{0})|^{p}\leq  \frac{\epsilon p}{2(p-1)}F(w(t)), \quad \text{and}\quad\left(\frac{t_{0}}{R_{1}}\right)^{\frac{p(n-1)}{p-1}}-1<\frac{\epsilon}{2}.
    \end{equation*}
    Here, the first one holds as $F(w(t))\to +\infty$ as $t\to R_{1}$.

    Therefore we have 
    \begin{equation*}
        -w'(t)\leq (1+\epsilon)^{1/p} \left(\frac{p}{p-1}F(w(t))\right)^{1/p}.
    \end{equation*}
    Finally,
    \begin{equation*}
    \begin{split}
         \Psi_{H,p}(w(t))&=\left(\frac{p-1}{p}\right)^{1/p}\int_{w(t)}^{\infty}\frac{ds}{\left\{F(s)\right\}^{1/p}}\\
         %&=\left(\frac{p-1}{p}\right)^{1/p}\int_{t}^{R_{1}} \frac{w'(s)}{\{F(w(s))\}^{1/p}}\ ds\\
         &=\left(\frac{p-1}{p}\right)^{1/p}\int^{t}_{R_{1}} \frac{-w'(s)}{\{F(w(s))\}^{1/p}}\ ds\\
         &\leq (1+\epsilon)^{1/p}(t-R_{1}).
         \end{split}
    \end{equation*}
    From this, we obtain
    \begin{equation*}
            \Psi_{H,p}(v(x))=\Psi_{H,p}(w(H_{0}(x)))\leq (1+\epsilon)^{1/p}(H_{0}(x)-R_{1})
    \end{equation*}
    This gives 
    \begin{equation*}
        \limsup_{|x|\to R_{1}}\frac{\Psi_{H,p}(v(x))}{{H_{0}}(x)-R_{1}}\leq 1.\qedhere
    \end{equation*}
\end{proof}

  \begin{lemma}\label{lemma asym 2}
       Let $\Omega$ be a bounded domain in $\rn$ which has a uniform exterior ball condition with respect to the dual norm $H_{0}$. Let $u$ be a solution of \eqref{finsler large eqn}. Then  
       \begin{equation*}
           \limsup_{\delta_{H_{0}}(x)\to 0}\frac{\Psi_{H,p}\left(u(x)\right)}{\delta_{H_{0}}(x)}\leq 1.
       \end{equation*}
       
   \end{lemma} 
   \begin{proof}

   Let $x\in \Omega_{\mu}$, then there exists $z(x)\in \partial\Omega$ such that \eqref{rep of x in Omega mu} holds. Denoting $z=z(x)$, let $\mathcal{W}_{R}(x^{ext})$ be the exterior $H_{0}-$ball associated with $z$. 
   
   Consider the solution $v_{r}$ of \eqref{eqn annulus} on $A_{R-r}^{2R}(x^{ext})$ for some $0<r<R$. Denote $\Omega_{r}'=\Omega\cap A_{R-r}^{2R}(x^{ext})$. Both $u$ and $v_{r}$ solve $\Delta_{H}^{p}w=f(w)$ in $\Omega_{r}'$. Since $u\geq 0, v_{r}=0$ on $\partial\Omega_{r}'\cap\partial A_{R-r}^{2R}(x^{ext})$ and $v_{r}<\infty$ on $\partial\Omega_{r}'\cap\partial\Omega$, by the comparison principle $v_{r}\leq u$ on $\Omega'_{r}$. 
   
   Further, for $r_{1}\leq r_{2}$ we have $v_{r_{1}}(x)\geq v_{r_{2}}(x)$ for all $x\in A_{R-r_{1}}^{2R}(x^{ext})$ by the comparison principle. Then $v=\sup v_{r}\leq u$ on $\Omega'=\Omega\cap A^{2R}_{R}(x^{ext})$. Moreover, employing the arguments analogous to the proof of Theorem \ref{main theorem existence}, $v$ is a solution of \eqref{eqn annulus} on $A_{R}^{2R}(x^{ext})$. We also refer to \cite[Theorem 4.3]{DattaIndro2026} for the use of the monotonicity method in the case when the sequence of solutions, defined in nested domains, blows up at the boundary. 
    
       Thus, by \eqref{rep of x in Omega mu} and \eqref{rep of center of ball}, since $\Psi_{H,p}$ is decreasing, we obtain
       \begin{equation*}
           \frac{\Psi_{H,p}\left(u(x)\right)}{\delta_{H_{0}}(x)}=\frac{\Psi_{H,p}\left(u(x)\right)}{H_{0}(x-z)}\leq \frac{\Psi_{H,p}\left(v(x)\right)}{H_{0}(x-x^{ext})-R}
       \end{equation*}
       Now, given $\epsilon>0$, by Proposition \ref{prop annuls asym}, there exists $\delta>0$ such that whenever $H_{0}(x-x^{ext})-R\leq \delta$, one has
       \begin{equation*}
           \frac{\Psi_{H,p}\left(v(x)\right)}{H_{0}(x-x^{ext})-R}< 1+\epsilon.
       \end{equation*}
       Choosing $\delta<\mu$, where $\mu$ is defined before \eqref{rep of x in Omega mu}, we obtain for all $x\in \Omega_{\delta}$ that 
       \begin{equation*}
        \frac{\Psi_{H,p}\left(u(x)\right)}{\delta_{H_{0}}(x)}< 1+\epsilon
       \end{equation*}
       Which completes the proof.
   \end{proof}

Combining the above results, we obtain the following proof of Theorem \ref{dual asym}.
    \begin{proof}[\textbf{Proof of Theorem \ref{dual asym}}]
        Since $\Omega$ is $C^{2}$, it has a uniform interior and exterior ball condition with respect to $H_{0}$. Applying Lemma \ref{lemma asym 1} and Lemma \ref{lemma asym 2}, we get 
        \begin{equation*}
            1\leq \liminf_{\delta_{H_{0}}(x)\to 0}\frac{\Psi_{H,p}\left(u(x)\right)}{\delta_{H_{0}}(x)}\leq \lim_{\delta_{H_{0}}(x)\to 0}\frac{\Psi_{H,p}\left(u(x)\right)}{\delta_{H_{0}}(x)}\leq  \limsup_{\delta_{H_{0}}(x)\to 0}\frac{\Psi_{H,p}\left(u(x)\right)}{\delta_{H_{0}}(x)}\leq 1.\qedhere
        \end{equation*}
    \end{proof}

    \section{Uniqueness}\label{uniqueness}
    In this section, we assume that $f(t)=t^{q}$ and establish the uniqueness using the asymptotics presented in the previous section.

   
    Let $c(p,q)$ denote a constant whose value may change from line to line. From Example \ref{example f=tq}, we obtain
    \begin{equation*}
        \Phi_{H,p}(s):=\Psi^{-1}_{H,p}(s)=\frac{c(p,q)}{s^{\frac{p}{q+1-p}}}.
    \end{equation*}
   
    \begin{proof}[\textbf{Proof of Corollary  \ref{thrm asym for f=tq}}]
     Given $\alpha>0$,  by Theorem \ref{dual asym}, there exists $\delta>0$ such that for all $x\in \Omega_{\delta}$ one has 
    \begin{equation*}
        (1-\alpha)\delta_{H_{0}}(x)<\Psi_{H,p}\left(u(x)\right)<(1+\alpha)\delta_{H_{0}}(x).
    \end{equation*}
    Since $\Psi_{H,p}$ is decreasing , so is $\Phi_{H,p}$. This implies 
    \begin{equation*}
        (1+\alpha)^{\frac{-p}{q+1-p}}\Phi(\delta_{H_{0}}(x))<u(x)<(1-\alpha)^{\frac{-p}{q+1-p}}\Phi(\delta_{H_{0}}(x)).
    \end{equation*}
    Rewriting 
    \begin{equation*}
        (1+\alpha)^{\frac{-p}{q+1-p}}<\frac{u(x)}{\Phi(\delta_{H_{0}}(x))}<(1-\alpha)^{\frac{-p}{q+1-p}}.
    \end{equation*}
    Since $1-(1+\alpha)^{\frac{-p}{q+1-p}}\to 0$ and $(1-\alpha)^{\frac{-p}{q+1-p}}-1\to 0$ as $\alpha\to 0$, we infer \eqref{asym f=tq}.
    \end{proof}
   
   \begin{proof}[\textbf{Proof of Theorem \ref{thrm uniqueness}}]
       Let $u$ and $v$ be two solutions of \eqref{finsler large eqn} with $f(t)=t^{q}$, $q>p-1$. Then, 
    \begin{equation*}
        \lim_{\delta_{H_{0}}(x)\to 0}\frac{u(x)}{v(x)}=\lim_{\delta_{H_{0}}(x)\to 0}\frac{u(x)}{\Psi_{H,p}(\delta_{H_{0}}(x))}~\frac{\Psi_{H,p}(\delta_{H_{0}}(x))}{v(x)}=1.
    \end{equation*}
    Therefore, by Theorem \ref{comparison principle}, $u\leq v$ in $\Omega$. Interchanging the roles of $u$ and $v$, we obtain uniqueness.
   \end{proof} 
    
    \section{Appendix}
This appendix gathers basic existence results about the Finsler p-Laplacian, followed by a brief discussion about the regularity.

Let $\Omega$ be a bounded domain in $\rn$, consider the energy 
 \begin{equation*}
    J(u)=\int_{\Omega}\frac{1}{p}H^{p}\left(\D u\right)+F(u) \ dx, \quad \text{for all}\ u\in W^{1,p}(\Omega).
\end{equation*}
of 
\begin{equation}\label{dirchlet problem}
     \begin{cases}
        -\Delta_{H}^{p}u+f(u)=0\quad &\text{in }\Omega\\
        u=g &\text{on } \partial\Omega,
    \end{cases}
\end{equation}
for some $g\in W^{1,p}(\Omega)$ and $F$ is the primitive of $f$ as in \eqref{eqn primitive of f}. The positivity of $f$ plays a role in the existence of a solution, as given below. 
\begin{definition}\label{dfn weak solution to dirchlet prob}
    The function $u\in W^{1,p}(\Omega)$ is a weak solution of \eqref{dirchlet problem} if
    \begin{equation*}
        -\int_{\Omega}\left<\Fin{u}, \D \phi\right>\ dx=\int_{\Omega}f(u)\phi\ dx \quad \text{for all }\phi\in W_{0}^{1,p}(\Omega). 
    \end{equation*}
    and $u-g\in W_{0}^{1,p}(\Omega)$ with $f(u)\in L^{p'}(\Omega)$.
\end{definition}
\begin{remark}
    In the above definition, $u$ also satisfies the weak formulation for all $\Omega'\subset\subset\Omega$ with the test function $\phi\in W^{1,p}(\Omega')$
\end{remark}
Clearly, the minimiser of the functional $J$ on $X:=\{u~|~u-g\in W_{0}^{1,p}(\Omega) \text{ and }F(u)\in L^{1}(\Omega)\}$ is the solution of \eqref{dirchlet problem}. Next, we show that the minimisers exist:
\begin{lemma}
    The functional $J$ is coercive, bounded below and weakly lower semi-continuous.
\end{lemma}
\begin{proof}
    By Lemma \ref{lemma_properties of H}, since $F\geq 0$ we infer
    \begin{equation*}
        J(u)\geq \int_{\Omega}\frac{1}{p}H^{p}(\nabla u) \ dx\geq \frac{\theta_{1}^{p}}{p}\int_{\Omega}|\nabla u|^{p}\ dx\geq 0.
    \end{equation*}
    Thus, $J$ is coercive and bounded below. To show that $J$ is weakly lower semi-continuous, let us divide $J$ into two functionals and show that both are weakly lower semi-continuous. Since the functional $u\mapsto \int_{\Omega}H^{p}(\nabla u)\ dx$ is continuous and convex, by \cite[Corollary 3.9, Remark 6]{Brezis2011} it is weakly lower semi-continuous. For the other part, let $v_{n}\rightharpoonup v$ in $W^{1,p}(\Omega)$. If $p< n$, then by the compact Sobolev embedding $v_{n}\to v$ in $L^{p}(\Omega)$ and, if $p\geq n$, then choose $0<\alpha<n^{2}/(p-n)$, then $p<(n-\alpha)^{*}$ and $W^{1,p}(\Omega)\subset W^{1,n-\alpha}(\Omega)$. The compact embedding then gives $v_{n}\to v$ in $L^{p}(\Omega)$ follows. Thus, there is a subsequence $v_{n_{k}}$ such that $v_{n_{k}}\to v$ pointwise a.e. The continuity of $F$ implies $F(v_{n_{k}})\to F(v)$ pointwise  a.e. Finally, Fatou's lemma gives the weakly lower semi-continuity of $u\mapsto \int_{\Omega}F(u)\ dx$.
    \end{proof}
      In view of the above lemma
    \begin{proposition}
        The functional $J$ admits a unique minimiser in $X$.
    \end{proposition}
    \begin{proof}
    Let $v_{n}$ be a minimising sequence in $X$. By the weak compactness $v_{n}\rightharpoonup v$ in $W^{1,p}(\Omega)$ up to some subsequence still denoted $v_{n}$. Further, $v$ is a minimiser of $J$ in $W^{1,p}(\Omega)$. The space $W_{0}^{1,p}(\Omega)$ is weakly closed. Thus, $v-g\in W_{0}^{1,p}(\Omega)$. Finally, by Fatou's lemma $F(v)\in L^{1}(\Omega)$ implies $v\in X$. 

    Let $v_{1}$ and $v_{2}$ be two minimisers. 
    By the monotonicity assumption on $f$,
    \begin{equation*}
        \int_{\Omega}\left(f(v_{1})-f(v_{2})\right)(v_{1}-v_{2})\ dx \geq 0.
    \end{equation*}
    By Corollary \ref{strict convexity of H}
    \begin{equation*}
        \int_{\Omega}\left<\Fin{v_{1}}-\Fin{v_{2}}, \nabla (v_{1}-v_{2})\right> \ dx>c \int_{\Omega}H^{p}(v_{1}-v_{2})\ dx,
    \end{equation*}
    if $v_{1}\neq v_{2}$. But since $v_{1}-v_{2}\in W_{0}^{1,p}(\Omega)$
    \begin{equation*}
        \int_{\Omega}\left<\Fin{v_{1}}-\Fin{v_{2}}, \nabla(v_{1}-v_{2})\right> +\left(f(v_{1})-f(v_{2})\right)(v_{1}-v_{2})=0.
    \end{equation*}
    Combining all these we conclude, $v_{1}=v_{2}$.
    \end{proof}
  Finaly, we state the regularity 
  \begin{proposition}
      Let $u$ be a bounded solution of \eqref{dirchlet problem}, then $u\in C_{loc}^{1,\alpha}(\Omega)$ for some $\alpha\in (0,1)$.
  \end{proposition}
  \begin{proof}
      The strong convexity is equivalent to uniform ellipticity. Thus by \cite[Proposition 3.1]{CozziFarinaValdinoci2016}, \cite[Proposition 3.1]{CozziFarinaValdinoci2014} (see the proof of Proposition 3.1 \cite{CozziFarinaValdinoci2014}), and \cite{Tolksdorf1984} the solutions are in $C_{loc}^{1,\alpha}(\Omega)$ for some $\alpha\in (0,1)$.
  \end{proof}
    %{\color{red}If we drop $u\in C(\Omega)$ in comparison principle, we get almost everywhere, and using this in local bound will imply that $u$ is locally bounded. Also, the boundedness of the solution follows the maximum principle when $g$ is continuous. Maximum principle follows by our comparison principle}

    \section{Acknowledgements}
    The author is supported by PMRF grant (2302262).

    The author thanks Dr. Indranil Chowdhury (Department of Mathematics and Statistics,
Indian Institute of Technology Kanpur) for the support and valuable suggestions.
\renewcommand\refname{Bibliography}
\bibliographystyle{abbrv}
\bibliography{ref.bib}
\end{document}
